\documentclass[11pt]{article}
\usepackage[margin=1in]{geometry}
\usepackage[T1]{fontenc}
\usepackage[utf8]{inputenc}
\usepackage{lmodern}
\usepackage{booktabs}
\usepackage{multirow}
\usepackage{graphicx}
\usepackage{tabularx}
\usepackage{amsmath}
\usepackage{amssymb}
\usepackage{microtype}
\usepackage{pgfplots}
\usepgfplotslibrary{groupplots}
\pgfplotsset{compat=1.18}
\usepackage[hidelinks]{hyperref}
\usepackage[numbers,sort&compress]{natbib}

% generated by clef/eval/cascade_paper.py -- do not edit
\expandafter\def\csname num:chk:mtp0:e2:b8_mean\endcsname{5772}
\expandafter\def\csname num:chk:mtp0:e2:big1_agree\endcsname{94.4}
\expandafter\def\csname num:chk:mtp0:e2:big1v8_agree\endcsname{94.4}
\expandafter\def\csname num:chk:mtp0:e2:big8_agree\endcsname{93.8}
\expandafter\def\csname num:chk:mtp0:e2:big_alone_esc_mean\endcsname{1741}
\expandafter\def\csname num:chk:mtp0:e2:big_frozen_acc\endcsname{76.1}
\expandafter\def\csname num:chk:mtp0:e2:big_in_casc_mean\endcsname{1698}
\expandafter\def\csname num:chk:mtp0:e2:big_live_acc\endcsname{77.1}
\expandafter\def\csname num:chk:mtp0:e2:big_mean\endcsname{1736}
\expandafter\def\csname num:chk:mtp0:e2:big_p50\endcsname{1588}
\expandafter\def\csname num:chk:mtp0:e2:big_p95\endcsname{2080}
\expandafter\def\csname num:chk:mtp0:e2:c8_mean\endcsname{5720}
\expandafter\def\csname num:chk:mtp0:e2:comp_mean\endcsname{1987}
\expandafter\def\csname num:chk:mtp0:e2:comp_p50\endcsname{1840}
\expandafter\def\csname num:chk:mtp0:e2:comp_p95\endcsname{2615}
\expandafter\def\csname num:chk:mtp0:e2:complive_mean\endcsname{1987}
\expandafter\def\csname num:chk:mtp0:e2:complive_p50\endcsname{1840}
\expandafter\def\csname num:chk:mtp0:e2:complive_p95\endcsname{2615}
\expandafter\def\csname num:chk:mtp0:e2:esc_agree\endcsname{100.0}
\expandafter\def\csname num:chk:mtp0:e2:live8_acc\endcsname{76.5}
\expandafter\def\csname num:chk:mtp0:e2:live8_esc\endcsname{98.0}
\expandafter\def\csname num:chk:mtp0:e2:live_acc\endcsname{75.1}
\expandafter\def\csname num:chk:mtp0:e2:live_ci\endcsname{[-3.6, -0.6]}
\expandafter\def\csname num:chk:mtp0:e2:live_d_big\endcsname{-2.0}
\expandafter\def\csname num:chk:mtp0:e2:live_esc\endcsname{98.0}
\expandafter\def\csname num:chk:mtp0:e2:meas_mean\endcsname{1863}
\expandafter\def\csname num:chk:mtp0:e2:meas_p50\endcsname{1885}
\expandafter\def\csname num:chk:mtp0:e2:meas_p95\endcsname{2330}
\expandafter\def\csname num:chk:mtp0:e2:n\endcsname{497}
\expandafter\def\csname num:chk:mtp0:e2:naive_mean\endcsname{5080}
\expandafter\def\csname num:chk:mtp0:e2:naive_p50\endcsname{5075}
\expandafter\def\csname num:chk:mtp0:e2:naive_p95\endcsname{6702}
\expandafter\def\csname num:chk:mtp0:e2:overhead_ms\endcsname{-124}
\expandafter\def\csname num:chk:mtp0:e2:pred_agree\endcsname{94.8}
\expandafter\def\csname num:chk:mtp0:e2:router_ms\endcsname{0.05}
\expandafter\def\csname num:chk:mtp0:e2:sim_acc\endcsname{74.8}
\expandafter\def\csname num:chk:mtp0:e2:sim_esc\endcsname{98.0}
\expandafter\def\csname num:chk:mtp0:e2:small_agree\endcsname{100.0}
\expandafter\def\csname num:chk:mtp0:e2:small_dconf\endcsname{0.00}
\expandafter\def\csname num:chk:mtp0:e2:small_frozen_acc\endcsname{19.7}
\expandafter\def\csname num:chk:mtp0:e2:small_in_casc_mean\endcsname{199}
\expandafter\def\csname num:chk:mtp0:e2:small_live_acc\endcsname{19.7}
\expandafter\def\csname num:chk:mtp0:e2:small_mean\endcsname{281}
\expandafter\def\csname num:chk:mtp0:e2:small_p50\endcsname{213}
\expandafter\def\csname num:chk:mtp0:e2:small_p95\endcsname{755}
\expandafter\def\csname num:chk:mtp0:e2:speedup_meas\endcsname{0.9}
\expandafter\def\csname num:chk:mtp0:e2:speedup_pred\endcsname{0.9}
\expandafter\def\csname num:chk:mtp0:e2:tput_gain\endcsname{1.0}
\expandafter\def\csname num:chk:mtp0:e2:x_big8\endcsname{1.38}
\expandafter\def\csname num:chk:mtp0:e2:x_bound\endcsname{1.41}
\expandafter\def\csname num:chk:mtp0:e2:x_meas\endcsname{1.40}
\expandafter\def\csname num:chk:mtp0:e2:x_model\endcsname{1.38}
\expandafter\def\csname num:chk:mtp0:e2:x_naive\endcsname{4.03}
\expandafter\def\csname num:chk:mtp0:e2:x_small8\endcsname{5.09}
\expandafter\def\csname num:chk:mtp0:e3:b8_mean\endcsname{2394}
\expandafter\def\csname num:chk:mtp0:e3:big1_agree\endcsname{99.0}
\expandafter\def\csname num:chk:mtp0:e3:big1_rep_agree\endcsname{99.8}
\expandafter\def\csname num:chk:mtp0:e3:big1_rep_mean\endcsname{834}
\expandafter\def\csname num:chk:mtp0:e3:big1v8_agree\endcsname{99.5}
\expandafter\def\csname num:chk:mtp0:e3:big8_agree\endcsname{99.2}
\expandafter\def\csname num:chk:mtp0:e3:big_alone_esc_mean\endcsname{600}
\expandafter\def\csname num:chk:mtp0:e3:big_frozen_acc\endcsname{84.7}
\expandafter\def\csname num:chk:mtp0:e3:big_in_casc_mean\endcsname{918}
\expandafter\def\csname num:chk:mtp0:e3:big_live_acc\endcsname{84.6}
\expandafter\def\csname num:chk:mtp0:e3:big_live_missed\endcsname{5}
\expandafter\def\csname num:chk:mtp0:e3:big_mean\endcsname{605}
\expandafter\def\csname num:chk:mtp0:e3:big_p50\endcsname{589}
\expandafter\def\csname num:chk:mtp0:e3:big_p95\endcsname{755}
\expandafter\def\csname num:chk:mtp0:e3:c8_mean\endcsname{878}
\expandafter\def\csname num:chk:mtp0:e3:comp_mean\endcsname{290}
\expandafter\def\csname num:chk:mtp0:e3:comp_p50\endcsname{83}
\expandafter\def\csname num:chk:mtp0:e3:comp_p95\endcsname{794}
\expandafter\def\csname num:chk:mtp0:e3:complive_mean\endcsname{290}
\expandafter\def\csname num:chk:mtp0:e3:complive_p50\endcsname{83}
\expandafter\def\csname num:chk:mtp0:e3:complive_p95\endcsname{794}
\expandafter\def\csname num:chk:mtp0:e3:esc_agree\endcsname{100.0}
\expandafter\def\csname num:chk:mtp0:e3:live8_acc\endcsname{84.3}
\expandafter\def\csname num:chk:mtp0:e3:live8_esc\endcsname{33.1}
\expandafter\def\csname num:chk:mtp0:e3:live_acc\endcsname{84.6}
\expandafter\def\csname num:chk:mtp0:e3:live_ci\endcsname{[-0.3, +0.3]}
\expandafter\def\csname num:chk:mtp0:e3:live_d_big\endcsname{0.0}
\expandafter\def\csname num:chk:mtp0:e3:live_esc\endcsname{33.1}
\expandafter\def\csname num:chk:mtp0:e3:live_missed\endcsname{6}
\expandafter\def\csname num:chk:mtp0:e3:meas_mean\endcsname{383}
\expandafter\def\csname num:chk:mtp0:e3:meas_p50\endcsname{81}
\expandafter\def\csname num:chk:mtp0:e3:meas_p95\endcsname{886}
\expandafter\def\csname num:chk:mtp0:e3:n\endcsname{864}
\expandafter\def\csname num:chk:mtp0:e3:naive_mean\endcsname{746}
\expandafter\def\csname num:chk:mtp0:e3:naive_p50\endcsname{85}
\expandafter\def\csname num:chk:mtp0:e3:naive_p95\endcsname{2494}
\expandafter\def\csname num:chk:mtp0:e3:overhead_ms\endcsname{94}
\expandafter\def\csname num:chk:mtp0:e3:pred_agree\endcsname{99.2}
\expandafter\def\csname num:chk:mtp0:e3:rep_slowdown\endcsname{38}
\expandafter\def\csname num:chk:mtp0:e3:router_ms\endcsname{0.02}
\expandafter\def\csname num:chk:mtp0:e3:sim_acc\endcsname{84.7}
\expandafter\def\csname num:chk:mtp0:e3:sim_esc\endcsname{33.1}
\expandafter\def\csname num:chk:mtp0:e3:sim_missed\endcsname{6}
\expandafter\def\csname num:chk:mtp0:e3:small_agree\endcsname{100.0}
\expandafter\def\csname num:chk:mtp0:e3:small_dconf\endcsname{0.00}
\expandafter\def\csname num:chk:mtp0:e3:small_frozen_acc\endcsname{79.2}
\expandafter\def\csname num:chk:mtp0:e3:small_in_casc_mean\endcsname{79}
\expandafter\def\csname num:chk:mtp0:e3:small_live_acc\endcsname{79.2}
\expandafter\def\csname num:chk:mtp0:e3:small_mean\endcsname{91}
\expandafter\def\csname num:chk:mtp0:e3:small_p50\endcsname{77}
\expandafter\def\csname num:chk:mtp0:e3:small_p95\endcsname{226}
\expandafter\def\csname num:chk:mtp0:e3:speedup_meas\endcsname{1.6}
\expandafter\def\csname num:chk:mtp0:e3:speedup_pred\endcsname{2.1}
\expandafter\def\csname num:chk:mtp0:e3:tput_gain\endcsname{2.7}
\expandafter\def\csname num:chk:mtp0:e3:x_big8\endcsname{3.34}
\expandafter\def\csname num:chk:mtp0:e3:x_bound\endcsname{10.08}
\expandafter\def\csname num:chk:mtp0:e3:x_meas\endcsname{9.05}
\expandafter\def\csname num:chk:mtp0:e3:x_model\endcsname{9.89}
\expandafter\def\csname num:chk:mtp0:e3:x_naive\endcsname{27.63}
\expandafter\def\csname num:chk:mtp0:e3:x_small8\endcsname{13.44}
\expandafter\def\csname num:chk:mtp0:e4ent:b8_mean\endcsname{3297}
\expandafter\def\csname num:chk:mtp0:e4ent:big1_agree\endcsname{100.0}
\expandafter\def\csname num:chk:mtp0:e4ent:big1v8_agree\endcsname{97.8}
\expandafter\def\csname num:chk:mtp0:e4ent:big8_agree\endcsname{97.8}
\expandafter\def\csname num:chk:mtp0:e4ent:big_frozen_acc\endcsname{87.0}
\expandafter\def\csname num:chk:mtp0:e4ent:big_live_acc\endcsname{87.0}
\expandafter\def\csname num:chk:mtp0:e4ent:big_mean\endcsname{792}
\expandafter\def\csname num:chk:mtp0:e4ent:big_p50\endcsname{719}
\expandafter\def\csname num:chk:mtp0:e4ent:big_p95\endcsname{1378}
\expandafter\def\csname num:chk:mtp0:e4ent:c8_mean\endcsname{984}
\expandafter\def\csname num:chk:mtp0:e4ent:comp_mean\endcsname{198}
\expandafter\def\csname num:chk:mtp0:e4ent:comp_p50\endcsname{167}
\expandafter\def\csname num:chk:mtp0:e4ent:comp_p95\endcsname{538}
\expandafter\def\csname num:chk:mtp0:e4ent:complive_mean\endcsname{198}
\expandafter\def\csname num:chk:mtp0:e4ent:complive_p50\endcsname{167}
\expandafter\def\csname num:chk:mtp0:e4ent:complive_p95\endcsname{538}
\expandafter\def\csname num:chk:mtp0:e4ent:esc_agree\endcsname{100.0}
\expandafter\def\csname num:chk:mtp0:e4ent:live8_acc\endcsname{89.1}
\expandafter\def\csname num:chk:mtp0:e4ent:live8_esc\endcsname{0.0}
\expandafter\def\csname num:chk:mtp0:e4ent:live_acc\endcsname{89.1}
\expandafter\def\csname num:chk:mtp0:e4ent:live_ci\endcsname{[-2.2, +6.5]}
\expandafter\def\csname num:chk:mtp0:e4ent:live_d_big\endcsname{2.2}
\expandafter\def\csname num:chk:mtp0:e4ent:live_esc\endcsname{0.0}
\expandafter\def\csname num:chk:mtp0:e4ent:meas_mean\endcsname{131}
\expandafter\def\csname num:chk:mtp0:e4ent:meas_p50\endcsname{128}
\expandafter\def\csname num:chk:mtp0:e4ent:meas_p95\endcsname{197}
\expandafter\def\csname num:chk:mtp0:e4ent:n\endcsname{92}
\expandafter\def\csname num:chk:mtp0:e4ent:naive_mean\endcsname{145}
\expandafter\def\csname num:chk:mtp0:e4ent:naive_p50\endcsname{135}
\expandafter\def\csname num:chk:mtp0:e4ent:naive_p95\endcsname{210}
\expandafter\def\csname num:chk:mtp0:e4ent:overhead_ms\endcsname{-67}
\expandafter\def\csname num:chk:mtp0:e4ent:pred_agree\endcsname{100.0}
\expandafter\def\csname num:chk:mtp0:e4ent:router_ms\endcsname{0.02}
\expandafter\def\csname num:chk:mtp0:e4ent:sim_acc\endcsname{89.1}
\expandafter\def\csname num:chk:mtp0:e4ent:sim_esc\endcsname{0.0}
\expandafter\def\csname num:chk:mtp0:e4ent:small_agree\endcsname{100.0}
\expandafter\def\csname num:chk:mtp0:e4ent:small_dconf\endcsname{0.00}
\expandafter\def\csname num:chk:mtp0:e4ent:small_frozen_acc\endcsname{89.1}
\expandafter\def\csname num:chk:mtp0:e4ent:small_in_casc_mean\endcsname{131}
\expandafter\def\csname num:chk:mtp0:e4ent:small_live_acc\endcsname{89.1}
\expandafter\def\csname num:chk:mtp0:e4ent:small_mean\endcsname{198}
\expandafter\def\csname num:chk:mtp0:e4ent:small_p50\endcsname{167}
\expandafter\def\csname num:chk:mtp0:e4ent:small_p95\endcsname{538}
\expandafter\def\csname num:chk:mtp0:e4ent:speedup_meas\endcsname{6.1}
\expandafter\def\csname num:chk:mtp0:e4ent:speedup_pred\endcsname{4.0}
\expandafter\def\csname num:chk:mtp0:e4ent:tput_gain\endcsname{3.3}
\expandafter\def\csname num:chk:mtp0:e4ent:x_big8\endcsname{2.40}
\expandafter\def\csname num:chk:mtp0:e4ent:x_bound\endcsname{7.92}
\expandafter\def\csname num:chk:mtp0:e4ent:x_meas\endcsname{7.95}
\expandafter\def\csname num:chk:mtp0:e4ent:x_model\endcsname{8.10}
\expandafter\def\csname num:chk:mtp0:e4ent:x_naive\endcsname{40.49}
\expandafter\def\csname num:chk:mtp0:e4ent:x_small8\endcsname{7.92}
\expandafter\def\csname num:chk:mtp0:e4facts:b8_mean\endcsname{2416}
\expandafter\def\csname num:chk:mtp0:e4facts:big1_agree\endcsname{96.6}
\expandafter\def\csname num:chk:mtp0:e4facts:big1v8_agree\endcsname{97.5}
\expandafter\def\csname num:chk:mtp0:e4facts:big8_agree\endcsname{97.1}
\expandafter\def\csname num:chk:mtp0:e4facts:big_alone_esc_mean\endcsname{571}
\expandafter\def\csname num:chk:mtp0:e4facts:big_frozen_acc\endcsname{88.3}
\expandafter\def\csname num:chk:mtp0:e4facts:big_in_casc_mean\endcsname{668}
\expandafter\def\csname num:chk:mtp0:e4facts:big_live_acc\endcsname{87.9}
\expandafter\def\csname num:chk:mtp0:e4facts:big_mean\endcsname{570}
\expandafter\def\csname num:chk:mtp0:e4facts:big_p50\endcsname{567}
\expandafter\def\csname num:chk:mtp0:e4facts:big_p95\endcsname{594}
\expandafter\def\csname num:chk:mtp0:e4facts:c8_mean\endcsname{975}
\expandafter\def\csname num:chk:mtp0:e4facts:comp_mean\endcsname{321}
\expandafter\def\csname num:chk:mtp0:e4facts:comp_p50\endcsname{96}
\expandafter\def\csname num:chk:mtp0:e4facts:comp_p95\endcsname{680}
\expandafter\def\csname num:chk:mtp0:e4facts:complive_mean\endcsname{321}
\expandafter\def\csname num:chk:mtp0:e4facts:complive_p50\endcsname{96}
\expandafter\def\csname num:chk:mtp0:e4facts:complive_p95\endcsname{680}
\expandafter\def\csname num:chk:mtp0:e4facts:esc_agree\endcsname{100.0}
\expandafter\def\csname num:chk:mtp0:e4facts:live8_acc\endcsname{88.1}
\expandafter\def\csname num:chk:mtp0:e4facts:live8_esc\endcsname{39.9}
\expandafter\def\csname num:chk:mtp0:e4facts:live_acc\endcsname{87.9}
\expandafter\def\csname num:chk:mtp0:e4facts:live_ci\endcsname{[-0.8, +0.8]}
\expandafter\def\csname num:chk:mtp0:e4facts:live_d_big\endcsname{0.0}
\expandafter\def\csname num:chk:mtp0:e4facts:live_esc\endcsname{39.9}
\expandafter\def\csname num:chk:mtp0:e4facts:meas_mean\endcsname{357}
\expandafter\def\csname num:chk:mtp0:e4facts:meas_p50\endcsname{96}
\expandafter\def\csname num:chk:mtp0:e4facts:meas_p95\endcsname{783}
\expandafter\def\csname num:chk:mtp0:e4facts:n\endcsname{1691}
\expandafter\def\csname num:chk:mtp0:e4facts:naive_mean\endcsname{902}
\expandafter\def\csname num:chk:mtp0:e4facts:naive_p50\endcsname{100}
\expandafter\def\csname num:chk:mtp0:e4facts:naive_p95\endcsname{2373}
\expandafter\def\csname num:chk:mtp0:e4facts:overhead_ms\endcsname{36}
\expandafter\def\csname num:chk:mtp0:e4facts:pred_agree\endcsname{97.6}
\expandafter\def\csname num:chk:mtp0:e4facts:router_ms\endcsname{0.03}
\expandafter\def\csname num:chk:mtp0:e4facts:sim_acc\endcsname{88.2}
\expandafter\def\csname num:chk:mtp0:e4facts:sim_esc\endcsname{39.9}
\expandafter\def\csname num:chk:mtp0:e4facts:small_agree\endcsname{100.0}
\expandafter\def\csname num:chk:mtp0:e4facts:small_dconf\endcsname{0.00}
\expandafter\def\csname num:chk:mtp0:e4facts:small_frozen_acc\endcsname{83.6}
\expandafter\def\csname num:chk:mtp0:e4facts:small_in_casc_mean\endcsname{91}
\expandafter\def\csname num:chk:mtp0:e4facts:small_live_acc\endcsname{83.6}
\expandafter\def\csname num:chk:mtp0:e4facts:small_mean\endcsname{93}
\expandafter\def\csname num:chk:mtp0:e4facts:small_p50\endcsname{91}
\expandafter\def\csname num:chk:mtp0:e4facts:small_p95\endcsname{101}
\expandafter\def\csname num:chk:mtp0:e4facts:speedup_meas\endcsname{1.6}
\expandafter\def\csname num:chk:mtp0:e4facts:speedup_pred\endcsname{1.8}
\expandafter\def\csname num:chk:mtp0:e4facts:tput_gain\endcsname{2.5}
\expandafter\def\csname num:chk:mtp0:e4facts:x_big8\endcsname{3.31}
\expandafter\def\csname num:chk:mtp0:e4facts:x_bound\endcsname{8.30}
\expandafter\def\csname num:chk:mtp0:e4facts:x_meas\endcsname{8.19}
\expandafter\def\csname num:chk:mtp0:e4facts:x_model\endcsname{8.16}
\expandafter\def\csname num:chk:mtp0:e4facts:x_naive\endcsname{24.95}
\expandafter\def\csname num:chk:mtp0:e4facts:x_small8\endcsname{11.67}
\expandafter\def\csname num:chk:mtp0apc0:e2:b8_mean\endcsname{5039}
\expandafter\def\csname num:chk:mtp0apc0:e2:big1_agree\endcsname{93.4}
\expandafter\def\csname num:chk:mtp0apc0:e2:big1v8_agree\endcsname{94.6}
\expandafter\def\csname num:chk:mtp0apc0:e2:big8_agree\endcsname{95.8}
\expandafter\def\csname num:chk:mtp0apc0:e2:big_alone_esc_mean\endcsname{1392}
\expandafter\def\csname num:chk:mtp0apc0:e2:big_frozen_acc\endcsname{76.1}
\expandafter\def\csname num:chk:mtp0apc0:e2:big_in_casc_mean\endcsname{1344}
\expandafter\def\csname num:chk:mtp0apc0:e2:big_live_acc\endcsname{76.7}
\expandafter\def\csname num:chk:mtp0apc0:e2:big_mean\endcsname{1392}
\expandafter\def\csname num:chk:mtp0apc0:e2:big_p50\endcsname{1399}
\expandafter\def\csname num:chk:mtp0apc0:e2:big_p95\endcsname{1635}
\expandafter\def\csname num:chk:mtp0apc0:e2:c8_mean\endcsname{4951}
\expandafter\def\csname num:chk:mtp0apc0:e2:comp_mean\endcsname{1644}
\expandafter\def\csname num:chk:mtp0apc0:e2:comp_p50\endcsname{1631}
\expandafter\def\csname num:chk:mtp0apc0:e2:comp_p95\endcsname{2267}
\expandafter\def\csname num:chk:mtp0apc0:e2:complive_mean\endcsname{1644}
\expandafter\def\csname num:chk:mtp0apc0:e2:complive_p50\endcsname{1631}
\expandafter\def\csname num:chk:mtp0apc0:e2:complive_p95\endcsname{2267}
\expandafter\def\csname num:chk:mtp0apc0:e2:esc_agree\endcsname{100.0}
\expandafter\def\csname num:chk:mtp0apc0:e2:live8_acc\endcsname{73.8}
\expandafter\def\csname num:chk:mtp0apc0:e2:live8_esc\endcsname{98.0}
\expandafter\def\csname num:chk:mtp0apc0:e2:live_acc\endcsname{75.5}
\expandafter\def\csname num:chk:mtp0apc0:e2:live_ci\endcsname{[-2.4, -0.4]}
\expandafter\def\csname num:chk:mtp0apc0:e2:live_d_big\endcsname{-1.2}
\expandafter\def\csname num:chk:mtp0apc0:e2:live_esc\endcsname{98.0}
\expandafter\def\csname num:chk:mtp0apc0:e2:meas_mean\endcsname{1517}
\expandafter\def\csname num:chk:mtp0apc0:e2:meas_p50\endcsname{1538}
\expandafter\def\csname num:chk:mtp0apc0:e2:meas_p95\endcsname{1855}
\expandafter\def\csname num:chk:mtp0apc0:e2:n\endcsname{497}
\expandafter\def\csname num:chk:mtp0apc0:e2:naive_mean\endcsname{5080}
\expandafter\def\csname num:chk:mtp0apc0:e2:naive_p50\endcsname{5075}
\expandafter\def\csname num:chk:mtp0apc0:e2:naive_p95\endcsname{6702}
\expandafter\def\csname num:chk:mtp0apc0:e2:overhead_ms\endcsname{-127}
\expandafter\def\csname num:chk:mtp0apc0:e2:pred_agree\endcsname{93.6}
\expandafter\def\csname num:chk:mtp0apc0:e2:router_ms\endcsname{0.05}
\expandafter\def\csname num:chk:mtp0apc0:e2:sim_acc\endcsname{74.8}
\expandafter\def\csname num:chk:mtp0apc0:e2:sim_esc\endcsname{98.0}
\expandafter\def\csname num:chk:mtp0apc0:e2:small_agree\endcsname{100.0}
\expandafter\def\csname num:chk:mtp0apc0:e2:small_dconf\endcsname{0.00}
\expandafter\def\csname num:chk:mtp0apc0:e2:small_frozen_acc\endcsname{19.7}
\expandafter\def\csname num:chk:mtp0apc0:e2:small_in_casc_mean\endcsname{200}
\expandafter\def\csname num:chk:mtp0apc0:e2:small_live_acc\endcsname{19.7}
\expandafter\def\csname num:chk:mtp0apc0:e2:small_mean\endcsname{281}
\expandafter\def\csname num:chk:mtp0apc0:e2:small_p50\endcsname{213}
\expandafter\def\csname num:chk:mtp0apc0:e2:small_p95\endcsname{755}
\expandafter\def\csname num:chk:mtp0apc0:e2:speedup_meas\endcsname{0.9}
\expandafter\def\csname num:chk:mtp0apc0:e2:speedup_pred\endcsname{0.8}
\expandafter\def\csname num:chk:mtp0apc0:e2:tput_gain\endcsname{1.0}
\expandafter\def\csname num:chk:mtp0apc0:e2:x_big8\endcsname{1.58}
\expandafter\def\csname num:chk:mtp0apc0:e2:x_bound\endcsname{1.62}
\expandafter\def\csname num:chk:mtp0apc0:e2:x_meas\endcsname{1.61}
\expandafter\def\csname num:chk:mtp0apc0:e2:x_model\endcsname{1.58}
\expandafter\def\csname num:chk:mtp0apc0:e2:x_naive\endcsname{4.86}
\expandafter\def\csname num:chk:mtp0apc0:e2:x_small8\endcsname{5.09}
\expandafter\def\csname num:chk:mtp0apc0:e3:b8_mean\endcsname{2097}
\expandafter\def\csname num:chk:mtp0apc0:e3:big1_agree\endcsname{98.8}
\expandafter\def\csname num:chk:mtp0apc0:e3:big1_rep_agree\endcsname{100.0}
\expandafter\def\csname num:chk:mtp0apc0:e3:big1_rep_mean\endcsname{598}
\expandafter\def\csname num:chk:mtp0apc0:e3:big1v8_agree\endcsname{99.0}
\expandafter\def\csname num:chk:mtp0apc0:e3:big8_agree\endcsname{99.4}
\expandafter\def\csname num:chk:mtp0apc0:e3:big_alone_esc_mean\endcsname{608}
\expandafter\def\csname num:chk:mtp0apc0:e3:big_frozen_acc\endcsname{84.7}
\expandafter\def\csname num:chk:mtp0apc0:e3:big_in_casc_mean\endcsname{575}
\expandafter\def\csname num:chk:mtp0apc0:e3:big_live_acc\endcsname{84.0}
\expandafter\def\csname num:chk:mtp0apc0:e3:big_live_missed\endcsname{6}
\expandafter\def\csname num:chk:mtp0apc0:e3:big_mean\endcsname{611}
\expandafter\def\csname num:chk:mtp0apc0:e3:big_p50\endcsname{595}
\expandafter\def\csname num:chk:mtp0apc0:e3:big_p95\endcsname{758}
\expandafter\def\csname num:chk:mtp0apc0:e3:c8_mean\endcsname{835}
\expandafter\def\csname num:chk:mtp0apc0:e3:comp_mean\endcsname{292}
\expandafter\def\csname num:chk:mtp0apc0:e3:comp_p50\endcsname{83}
\expandafter\def\csname num:chk:mtp0apc0:e3:comp_p95\endcsname{809}
\expandafter\def\csname num:chk:mtp0apc0:e3:complive_mean\endcsname{292}
\expandafter\def\csname num:chk:mtp0apc0:e3:complive_p50\endcsname{83}
\expandafter\def\csname num:chk:mtp0apc0:e3:complive_p95\endcsname{809}
\expandafter\def\csname num:chk:mtp0apc0:e3:esc_agree\endcsname{100.0}
\expandafter\def\csname num:chk:mtp0apc0:e3:live8_acc\endcsname{84.3}
\expandafter\def\csname num:chk:mtp0apc0:e3:live8_esc\endcsname{33.1}
\expandafter\def\csname num:chk:mtp0apc0:e3:live_acc\endcsname{84.1}
\expandafter\def\csname num:chk:mtp0apc0:e3:live_ci\endcsname{[+0.0, +0.3]}
\expandafter\def\csname num:chk:mtp0apc0:e3:live_d_big\endcsname{0.1}
\expandafter\def\csname num:chk:mtp0apc0:e3:live_esc\endcsname{33.1}
\expandafter\def\csname num:chk:mtp0apc0:e3:live_missed\endcsname{6}
\expandafter\def\csname num:chk:mtp0apc0:e3:meas_mean\endcsname{270}
\expandafter\def\csname num:chk:mtp0apc0:e3:meas_p50\endcsname{82}
\expandafter\def\csname num:chk:mtp0apc0:e3:meas_p95\endcsname{745}
\expandafter\def\csname num:chk:mtp0apc0:e3:n\endcsname{864}
\expandafter\def\csname num:chk:mtp0apc0:e3:naive_mean\endcsname{746}
\expandafter\def\csname num:chk:mtp0apc0:e3:naive_p50\endcsname{85}
\expandafter\def\csname num:chk:mtp0apc0:e3:naive_p95\endcsname{2494}
\expandafter\def\csname num:chk:mtp0apc0:e3:overhead_ms\endcsname{-22}
\expandafter\def\csname num:chk:mtp0apc0:e3:pred_agree\endcsname{99.2}
\expandafter\def\csname num:chk:mtp0apc0:e3:rep_slowdown\endcsname{-2}
\expandafter\def\csname num:chk:mtp0apc0:e3:router_ms\endcsname{0.02}
\expandafter\def\csname num:chk:mtp0apc0:e3:sim_acc\endcsname{84.7}
\expandafter\def\csname num:chk:mtp0apc0:e3:sim_esc\endcsname{33.1}
\expandafter\def\csname num:chk:mtp0apc0:e3:sim_missed\endcsname{6}
\expandafter\def\csname num:chk:mtp0apc0:e3:small_agree\endcsname{100.0}
\expandafter\def\csname num:chk:mtp0apc0:e3:small_dconf\endcsname{0.00}
\expandafter\def\csname num:chk:mtp0apc0:e3:small_frozen_acc\endcsname{79.2}
\expandafter\def\csname num:chk:mtp0apc0:e3:small_in_casc_mean\endcsname{80}
\expandafter\def\csname num:chk:mtp0apc0:e3:small_live_acc\endcsname{79.2}
\expandafter\def\csname num:chk:mtp0apc0:e3:small_mean\endcsname{91}
\expandafter\def\csname num:chk:mtp0apc0:e3:small_p50\endcsname{77}
\expandafter\def\csname num:chk:mtp0apc0:e3:small_p95\endcsname{226}
\expandafter\def\csname num:chk:mtp0apc0:e3:speedup_meas\endcsname{2.3}
\expandafter\def\csname num:chk:mtp0apc0:e3:speedup_pred\endcsname{2.1}
\expandafter\def\csname num:chk:mtp0apc0:e3:tput_gain\endcsname{2.5}
\expandafter\def\csname num:chk:mtp0apc0:e3:x_big8\endcsname{3.81}
\expandafter\def\csname num:chk:mtp0apc0:e3:x_bound\endcsname{11.50}
\expandafter\def\csname num:chk:mtp0apc0:e3:x_meas\endcsname{9.51}
\expandafter\def\csname num:chk:mtp0apc0:e3:x_model\endcsname{11.19}
\expandafter\def\csname num:chk:mtp0apc0:e3:x_naive\endcsname{27.39}
\expandafter\def\csname num:chk:mtp0apc0:e3:x_small8\endcsname{13.44}
\expandafter\def\csname num:chk:mtp0apc0:e4ent:b8_mean\endcsname{2894}
\expandafter\def\csname num:chk:mtp0apc0:e4ent:big1_agree\endcsname{98.9}
\expandafter\def\csname num:chk:mtp0apc0:e4ent:big1v8_agree\endcsname{98.9}
\expandafter\def\csname num:chk:mtp0apc0:e4ent:big8_agree\endcsname{100.0}
\expandafter\def\csname num:chk:mtp0apc0:e4ent:big_frozen_acc\endcsname{87.0}
\expandafter\def\csname num:chk:mtp0apc0:e4ent:big_live_acc\endcsname{85.9}
\expandafter\def\csname num:chk:mtp0apc0:e4ent:big_mean\endcsname{748}
\expandafter\def\csname num:chk:mtp0apc0:e4ent:big_p50\endcsname{706}
\expandafter\def\csname num:chk:mtp0apc0:e4ent:big_p95\endcsname{1029}
\expandafter\def\csname num:chk:mtp0apc0:e4ent:c8_mean\endcsname{983}
\expandafter\def\csname num:chk:mtp0apc0:e4ent:comp_mean\endcsname{198}
\expandafter\def\csname num:chk:mtp0apc0:e4ent:comp_p50\endcsname{167}
\expandafter\def\csname num:chk:mtp0apc0:e4ent:comp_p95\endcsname{538}
\expandafter\def\csname num:chk:mtp0apc0:e4ent:complive_mean\endcsname{198}
\expandafter\def\csname num:chk:mtp0apc0:e4ent:complive_p50\endcsname{167}
\expandafter\def\csname num:chk:mtp0apc0:e4ent:complive_p95\endcsname{538}
\expandafter\def\csname num:chk:mtp0apc0:e4ent:esc_agree\endcsname{100.0}
\expandafter\def\csname num:chk:mtp0apc0:e4ent:live8_acc\endcsname{89.1}
\expandafter\def\csname num:chk:mtp0apc0:e4ent:live8_esc\endcsname{0.0}
\expandafter\def\csname num:chk:mtp0apc0:e4ent:live_acc\endcsname{89.1}
\expandafter\def\csname num:chk:mtp0apc0:e4ent:live_ci\endcsname{[-1.1, +7.6]}
\expandafter\def\csname num:chk:mtp0apc0:e4ent:live_d_big\endcsname{3.3}
\expandafter\def\csname num:chk:mtp0apc0:e4ent:live_esc\endcsname{0.0}
\expandafter\def\csname num:chk:mtp0apc0:e4ent:meas_mean\endcsname{131}
\expandafter\def\csname num:chk:mtp0apc0:e4ent:meas_p50\endcsname{129}
\expandafter\def\csname num:chk:mtp0apc0:e4ent:meas_p95\endcsname{198}
\expandafter\def\csname num:chk:mtp0apc0:e4ent:n\endcsname{92}
\expandafter\def\csname num:chk:mtp0apc0:e4ent:naive_mean\endcsname{145}
\expandafter\def\csname num:chk:mtp0apc0:e4ent:naive_p50\endcsname{135}
\expandafter\def\csname num:chk:mtp0apc0:e4ent:naive_p95\endcsname{210}
\expandafter\def\csname num:chk:mtp0apc0:e4ent:overhead_ms\endcsname{-66}
\expandafter\def\csname num:chk:mtp0apc0:e4ent:pred_agree\endcsname{100.0}
\expandafter\def\csname num:chk:mtp0apc0:e4ent:router_ms\endcsname{0.01}
\expandafter\def\csname num:chk:mtp0apc0:e4ent:sim_acc\endcsname{89.1}
\expandafter\def\csname num:chk:mtp0apc0:e4ent:sim_esc\endcsname{0.0}
\expandafter\def\csname num:chk:mtp0apc0:e4ent:small_agree\endcsname{100.0}
\expandafter\def\csname num:chk:mtp0apc0:e4ent:small_dconf\endcsname{0.00}
\expandafter\def\csname num:chk:mtp0apc0:e4ent:small_frozen_acc\endcsname{89.1}
\expandafter\def\csname num:chk:mtp0apc0:e4ent:small_in_casc_mean\endcsname{131}
\expandafter\def\csname num:chk:mtp0apc0:e4ent:small_live_acc\endcsname{89.1}
\expandafter\def\csname num:chk:mtp0apc0:e4ent:small_mean\endcsname{198}
\expandafter\def\csname num:chk:mtp0apc0:e4ent:small_p50\endcsname{167}
\expandafter\def\csname num:chk:mtp0apc0:e4ent:small_p95\endcsname{538}
\expandafter\def\csname num:chk:mtp0apc0:e4ent:speedup_meas\endcsname{5.7}
\expandafter\def\csname num:chk:mtp0apc0:e4ent:speedup_pred\endcsname{3.8}
\expandafter\def\csname num:chk:mtp0apc0:e4ent:tput_gain\endcsname{2.9}
\expandafter\def\csname num:chk:mtp0apc0:e4ent:x_big8\endcsname{2.71}
\expandafter\def\csname num:chk:mtp0apc0:e4ent:x_bound\endcsname{7.92}
\expandafter\def\csname num:chk:mtp0apc0:e4ent:x_meas\endcsname{7.95}
\expandafter\def\csname num:chk:mtp0apc0:e4ent:x_model\endcsname{8.10}
\expandafter\def\csname num:chk:mtp0apc0:e4ent:x_naive\endcsname{40.49}
\expandafter\def\csname num:chk:mtp0apc0:e4ent:x_small8\endcsname{7.92}
\expandafter\def\csname num:chk:mtp0apc0:e4facts:b8_mean\endcsname{1967}
\expandafter\def\csname num:chk:mtp0apc0:e4facts:big1_agree\endcsname{97.6}
\expandafter\def\csname num:chk:mtp0apc0:e4facts:big1v8_agree\endcsname{97.0}
\expandafter\def\csname num:chk:mtp0apc0:e4facts:big8_agree\endcsname{96.9}
\expandafter\def\csname num:chk:mtp0apc0:e4facts:big_alone_esc_mean\endcsname{574}
\expandafter\def\csname num:chk:mtp0apc0:e4facts:big_frozen_acc\endcsname{88.3}
\expandafter\def\csname num:chk:mtp0apc0:e4facts:big_in_casc_mean\endcsname{543}
\expandafter\def\csname num:chk:mtp0apc0:e4facts:big_live_acc\endcsname{88.3}
\expandafter\def\csname num:chk:mtp0apc0:e4facts:big_mean\endcsname{574}
\expandafter\def\csname num:chk:mtp0apc0:e4facts:big_p50\endcsname{570}
\expandafter\def\csname num:chk:mtp0apc0:e4facts:big_p95\endcsname{600}
\expandafter\def\csname num:chk:mtp0apc0:e4facts:c8_mean\endcsname{953}
\expandafter\def\csname num:chk:mtp0apc0:e4facts:comp_mean\endcsname{322}
\expandafter\def\csname num:chk:mtp0apc0:e4facts:comp_p50\endcsname{96}
\expandafter\def\csname num:chk:mtp0apc0:e4facts:comp_p95\endcsname{684}
\expandafter\def\csname num:chk:mtp0apc0:e4facts:complive_mean\endcsname{322}
\expandafter\def\csname num:chk:mtp0apc0:e4facts:complive_p50\endcsname{96}
\expandafter\def\csname num:chk:mtp0apc0:e4facts:complive_p95\endcsname{684}
\expandafter\def\csname num:chk:mtp0apc0:e4facts:esc_agree\endcsname{100.0}
\expandafter\def\csname num:chk:mtp0apc0:e4facts:live8_acc\endcsname{88.2}
\expandafter\def\csname num:chk:mtp0apc0:e4facts:live8_esc\endcsname{39.9}
\expandafter\def\csname num:chk:mtp0apc0:e4facts:live_acc\endcsname{88.2}
\expandafter\def\csname num:chk:mtp0apc0:e4facts:live_ci\endcsname{[-0.7, +0.4]}
\expandafter\def\csname num:chk:mtp0apc0:e4facts:live_d_big\endcsname{-0.1}
\expandafter\def\csname num:chk:mtp0apc0:e4facts:live_esc\endcsname{39.9}
\expandafter\def\csname num:chk:mtp0apc0:e4facts:meas_mean\endcsname{308}
\expandafter\def\csname num:chk:mtp0apc0:e4facts:meas_p50\endcsname{97}
\expandafter\def\csname num:chk:mtp0apc0:e4facts:meas_p95\endcsname{652}
\expandafter\def\csname num:chk:mtp0apc0:e4facts:n\endcsname{1691}
\expandafter\def\csname num:chk:mtp0apc0:e4facts:naive_mean\endcsname{902}
\expandafter\def\csname num:chk:mtp0apc0:e4facts:naive_p50\endcsname{100}
\expandafter\def\csname num:chk:mtp0apc0:e4facts:naive_p95\endcsname{2373}
\expandafter\def\csname num:chk:mtp0apc0:e4facts:overhead_ms\endcsname{-14}
\expandafter\def\csname num:chk:mtp0apc0:e4facts:pred_agree\endcsname{98.2}
\expandafter\def\csname num:chk:mtp0apc0:e4facts:router_ms\endcsname{0.03}
\expandafter\def\csname num:chk:mtp0apc0:e4facts:sim_acc\endcsname{88.2}
\expandafter\def\csname num:chk:mtp0apc0:e4facts:sim_esc\endcsname{39.9}
\expandafter\def\csname num:chk:mtp0apc0:e4facts:small_agree\endcsname{100.0}
\expandafter\def\csname num:chk:mtp0apc0:e4facts:small_dconf\endcsname{0.00}
\expandafter\def\csname num:chk:mtp0apc0:e4facts:small_frozen_acc\endcsname{83.6}
\expandafter\def\csname num:chk:mtp0apc0:e4facts:small_in_casc_mean\endcsname{91}
\expandafter\def\csname num:chk:mtp0apc0:e4facts:small_live_acc\endcsname{83.6}
\expandafter\def\csname num:chk:mtp0apc0:e4facts:small_mean\endcsname{93}
\expandafter\def\csname num:chk:mtp0apc0:e4facts:small_p50\endcsname{91}
\expandafter\def\csname num:chk:mtp0apc0:e4facts:small_p95\endcsname{101}
\expandafter\def\csname num:chk:mtp0apc0:e4facts:speedup_meas\endcsname{1.9}
\expandafter\def\csname num:chk:mtp0apc0:e4facts:speedup_pred\endcsname{1.8}
\expandafter\def\csname num:chk:mtp0apc0:e4facts:tput_gain\endcsname{2.1}
\expandafter\def\csname num:chk:mtp0apc0:e4facts:x_big8\endcsname{4.06}
\expandafter\def\csname num:chk:mtp0apc0:e4facts:x_bound\endcsname{10.20}
\expandafter\def\csname num:chk:mtp0apc0:e4facts:x_meas\endcsname{8.37}
\expandafter\def\csname num:chk:mtp0apc0:e4facts:x_model\endcsname{9.93}
\expandafter\def\csname num:chk:mtp0apc0:e4facts:x_naive\endcsname{24.85}
\expandafter\def\csname num:chk:mtp0apc0:e4facts:x_small8\endcsname{11.67}
\expandafter\def\csname num:chk:mtp2:e2:b8_mean\endcsname{4596}
\expandafter\def\csname num:chk:mtp2:e2:big1_agree\endcsname{93.2}
\expandafter\def\csname num:chk:mtp2:e2:big1v8_agree\endcsname{93.4}
\expandafter\def\csname num:chk:mtp2:e2:big8_agree\endcsname{94.8}
\expandafter\def\csname num:chk:mtp2:e2:big_alone_esc_mean\endcsname{1077}
\expandafter\def\csname num:chk:mtp2:e2:big_frozen_acc\endcsname{76.1}
\expandafter\def\csname num:chk:mtp2:e2:big_in_casc_mean\endcsname{1013}
\expandafter\def\csname num:chk:mtp2:e2:big_live_acc\endcsname{75.9}
\expandafter\def\csname num:chk:mtp2:e2:big_mean\endcsname{1077}
\expandafter\def\csname num:chk:mtp2:e2:big_p50\endcsname{1066}
\expandafter\def\csname num:chk:mtp2:e2:big_p95\endcsname{1289}
\expandafter\def\csname num:chk:mtp2:e2:c8_mean\endcsname{4633}
\expandafter\def\csname num:chk:mtp2:e2:comp_mean\endcsname{1336}
\expandafter\def\csname num:chk:mtp2:e2:comp_p50\endcsname{1293}
\expandafter\def\csname num:chk:mtp2:e2:comp_p95\endcsname{1908}
\expandafter\def\csname num:chk:mtp2:e2:complive_mean\endcsname{1336}
\expandafter\def\csname num:chk:mtp2:e2:complive_p50\endcsname{1293}
\expandafter\def\csname num:chk:mtp2:e2:complive_p95\endcsname{1908}
\expandafter\def\csname num:chk:mtp2:e2:esc_agree\endcsname{100.0}
\expandafter\def\csname num:chk:mtp2:e2:live8_acc\endcsname{75.1}
\expandafter\def\csname num:chk:mtp2:e2:live8_esc\endcsname{98.0}
\expandafter\def\csname num:chk:mtp2:e2:live_acc\endcsname{74.6}
\expandafter\def\csname num:chk:mtp2:e2:live_ci\endcsname{[-2.4, -0.4]}
\expandafter\def\csname num:chk:mtp2:e2:live_d_big\endcsname{-1.2}
\expandafter\def\csname num:chk:mtp2:e2:live_esc\endcsname{98.0}
\expandafter\def\csname num:chk:mtp2:e2:meas_mean\endcsname{1227}
\expandafter\def\csname num:chk:mtp2:e2:meas_p50\endcsname{1212}
\expandafter\def\csname num:chk:mtp2:e2:meas_p95\endcsname{1575}
\expandafter\def\csname num:chk:mtp2:e2:n\endcsname{497}
\expandafter\def\csname num:chk:mtp2:e2:naive_mean\endcsname{5080}
\expandafter\def\csname num:chk:mtp2:e2:naive_p50\endcsname{5075}
\expandafter\def\csname num:chk:mtp2:e2:naive_p95\endcsname{6702}
\expandafter\def\csname num:chk:mtp2:e2:overhead_ms\endcsname{-110}
\expandafter\def\csname num:chk:mtp2:e2:pred_agree\endcsname{93.2}
\expandafter\def\csname num:chk:mtp2:e2:router_ms\endcsname{0.05}
\expandafter\def\csname num:chk:mtp2:e2:sim_acc\endcsname{74.8}
\expandafter\def\csname num:chk:mtp2:e2:sim_esc\endcsname{98.0}
\expandafter\def\csname num:chk:mtp2:e2:small_agree\endcsname{100.0}
\expandafter\def\csname num:chk:mtp2:e2:small_dconf\endcsname{0.00}
\expandafter\def\csname num:chk:mtp2:e2:small_frozen_acc\endcsname{19.7}
\expandafter\def\csname num:chk:mtp2:e2:small_in_casc_mean\endcsname{233}
\expandafter\def\csname num:chk:mtp2:e2:small_live_acc\endcsname{19.7}
\expandafter\def\csname num:chk:mtp2:e2:small_mean\endcsname{281}
\expandafter\def\csname num:chk:mtp2:e2:small_p50\endcsname{213}
\expandafter\def\csname num:chk:mtp2:e2:small_p95\endcsname{755}
\expandafter\def\csname num:chk:mtp2:e2:speedup_meas\endcsname{0.9}
\expandafter\def\csname num:chk:mtp2:e2:speedup_pred\endcsname{0.8}
\expandafter\def\csname num:chk:mtp2:e2:tput_gain\endcsname{1.0}
\expandafter\def\csname num:chk:mtp2:e2:x_big8\endcsname{1.74}
\expandafter\def\csname num:chk:mtp2:e2:x_bound\endcsname{1.77}
\expandafter\def\csname num:chk:mtp2:e2:x_meas\endcsname{1.72}
\expandafter\def\csname num:chk:mtp2:e2:x_model\endcsname{1.75}
\expandafter\def\csname num:chk:mtp2:e2:x_naive\endcsname{5.99}
\expandafter\def\csname num:chk:mtp2:e2:x_small8\endcsname{5.09}
\expandafter\def\csname num:chk:mtp2:e3:b8_mean\endcsname{1960}
\expandafter\def\csname num:chk:mtp2:e3:big1_agree\endcsname{99.0}
\expandafter\def\csname num:chk:mtp2:e3:big1_rep_agree\endcsname{100.0}
\expandafter\def\csname num:chk:mtp2:e3:big1_rep_mean\endcsname{519}
\expandafter\def\csname num:chk:mtp2:e3:big1v8_agree\endcsname{99.1}
\expandafter\def\csname num:chk:mtp2:e3:big8_agree\endcsname{99.0}
\expandafter\def\csname num:chk:mtp2:e3:big_alone_esc_mean\endcsname{519}
\expandafter\def\csname num:chk:mtp2:e3:big_frozen_acc\endcsname{84.7}
\expandafter\def\csname num:chk:mtp2:e3:big_in_casc_mean\endcsname{484}
\expandafter\def\csname num:chk:mtp2:e3:big_live_acc\endcsname{84.4}
\expandafter\def\csname num:chk:mtp2:e3:big_live_missed\endcsname{5}
\expandafter\def\csname num:chk:mtp2:e3:big_mean\endcsname{528}
\expandafter\def\csname num:chk:mtp2:e3:big_p50\endcsname{519}
\expandafter\def\csname num:chk:mtp2:e3:big_p95\endcsname{618}
\expandafter\def\csname num:chk:mtp2:e3:c8_mean\endcsname{809}
\expandafter\def\csname num:chk:mtp2:e3:comp_mean\endcsname{263}
\expandafter\def\csname num:chk:mtp2:e3:comp_p50\endcsname{83}
\expandafter\def\csname num:chk:mtp2:e3:comp_p95\endcsname{660}
\expandafter\def\csname num:chk:mtp2:e3:complive_mean\endcsname{263}
\expandafter\def\csname num:chk:mtp2:e3:complive_p50\endcsname{83}
\expandafter\def\csname num:chk:mtp2:e3:complive_p95\endcsname{660}
\expandafter\def\csname num:chk:mtp2:e3:esc_agree\endcsname{100.0}
\expandafter\def\csname num:chk:mtp2:e3:live8_acc\endcsname{84.7}
\expandafter\def\csname num:chk:mtp2:e3:live8_esc\endcsname{33.1}
\expandafter\def\csname num:chk:mtp2:e3:live_acc\endcsname{84.5}
\expandafter\def\csname num:chk:mtp2:e3:live_ci\endcsname{[+0.0, +0.3]}
\expandafter\def\csname num:chk:mtp2:e3:live_d_big\endcsname{0.1}
\expandafter\def\csname num:chk:mtp2:e3:live_esc\endcsname{33.1}
\expandafter\def\csname num:chk:mtp2:e3:live_missed\endcsname{5}
\expandafter\def\csname num:chk:mtp2:e3:meas_mean\endcsname{241}
\expandafter\def\csname num:chk:mtp2:e3:meas_p50\endcsname{82}
\expandafter\def\csname num:chk:mtp2:e3:meas_p95\endcsname{620}
\expandafter\def\csname num:chk:mtp2:e3:n\endcsname{864}
\expandafter\def\csname num:chk:mtp2:e3:naive_mean\endcsname{746}
\expandafter\def\csname num:chk:mtp2:e3:naive_p50\endcsname{85}
\expandafter\def\csname num:chk:mtp2:e3:naive_p95\endcsname{2494}
\expandafter\def\csname num:chk:mtp2:e3:overhead_ms\endcsname{-22}
\expandafter\def\csname num:chk:mtp2:e3:pred_agree\endcsname{99.3}
\expandafter\def\csname num:chk:mtp2:e3:rep_slowdown\endcsname{-2}
\expandafter\def\csname num:chk:mtp2:e3:router_ms\endcsname{0.03}
\expandafter\def\csname num:chk:mtp2:e3:sim_acc\endcsname{84.7}
\expandafter\def\csname num:chk:mtp2:e3:sim_esc\endcsname{33.1}
\expandafter\def\csname num:chk:mtp2:e3:sim_missed\endcsname{6}
\expandafter\def\csname num:chk:mtp2:e3:small_agree\endcsname{100.0}
\expandafter\def\csname num:chk:mtp2:e3:small_dconf\endcsname{0.00}
\expandafter\def\csname num:chk:mtp2:e3:small_frozen_acc\endcsname{79.2}
\expandafter\def\csname num:chk:mtp2:e3:small_in_casc_mean\endcsname{81}
\expandafter\def\csname num:chk:mtp2:e3:small_live_acc\endcsname{79.2}
\expandafter\def\csname num:chk:mtp2:e3:small_mean\endcsname{91}
\expandafter\def\csname num:chk:mtp2:e3:small_p50\endcsname{77}
\expandafter\def\csname num:chk:mtp2:e3:small_p95\endcsname{226}
\expandafter\def\csname num:chk:mtp2:e3:speedup_meas\endcsname{2.2}
\expandafter\def\csname num:chk:mtp2:e3:speedup_pred\endcsname{2.0}
\expandafter\def\csname num:chk:mtp2:e3:tput_gain\endcsname{2.4}
\expandafter\def\csname num:chk:mtp2:e3:x_big8\endcsname{4.08}
\expandafter\def\csname num:chk:mtp2:e3:x_bound\endcsname{12.31}
\expandafter\def\csname num:chk:mtp2:e3:x_meas\endcsname{9.84}
\expandafter\def\csname num:chk:mtp2:e3:x_model\endcsname{11.86}
\expandafter\def\csname num:chk:mtp2:e3:x_naive\endcsname{30.44}
\expandafter\def\csname num:chk:mtp2:e3:x_small8\endcsname{13.44}
\expandafter\def\csname num:chk:mtp2:e4ent:b8_mean\endcsname{2965}
\expandafter\def\csname num:chk:mtp2:e4ent:big1_agree\endcsname{98.9}
\expandafter\def\csname num:chk:mtp2:e4ent:big1v8_agree\endcsname{97.8}
\expandafter\def\csname num:chk:mtp2:e4ent:big8_agree\endcsname{98.9}
\expandafter\def\csname num:chk:mtp2:e4ent:big_frozen_acc\endcsname{87.0}
\expandafter\def\csname num:chk:mtp2:e4ent:big_live_acc\endcsname{88.0}
\expandafter\def\csname num:chk:mtp2:e4ent:big_mean\endcsname{616}
\expandafter\def\csname num:chk:mtp2:e4ent:big_p50\endcsname{613}
\expandafter\def\csname num:chk:mtp2:e4ent:big_p95\endcsname{779}
\expandafter\def\csname num:chk:mtp2:e4ent:c8_mean\endcsname{981}
\expandafter\def\csname num:chk:mtp2:e4ent:comp_mean\endcsname{198}
\expandafter\def\csname num:chk:mtp2:e4ent:comp_p50\endcsname{167}
\expandafter\def\csname num:chk:mtp2:e4ent:comp_p95\endcsname{538}
\expandafter\def\csname num:chk:mtp2:e4ent:complive_mean\endcsname{198}
\expandafter\def\csname num:chk:mtp2:e4ent:complive_p50\endcsname{167}
\expandafter\def\csname num:chk:mtp2:e4ent:complive_p95\endcsname{538}
\expandafter\def\csname num:chk:mtp2:e4ent:esc_agree\endcsname{100.0}
\expandafter\def\csname num:chk:mtp2:e4ent:live8_acc\endcsname{89.1}
\expandafter\def\csname num:chk:mtp2:e4ent:live8_esc\endcsname{0.0}
\expandafter\def\csname num:chk:mtp2:e4ent:live_acc\endcsname{89.1}
\expandafter\def\csname num:chk:mtp2:e4ent:live_ci\endcsname{[-2.2, +5.4]}
\expandafter\def\csname num:chk:mtp2:e4ent:live_d_big\endcsname{1.1}
\expandafter\def\csname num:chk:mtp2:e4ent:live_esc\endcsname{0.0}
\expandafter\def\csname num:chk:mtp2:e4ent:meas_mean\endcsname{139}
\expandafter\def\csname num:chk:mtp2:e4ent:meas_p50\endcsname{132}
\expandafter\def\csname num:chk:mtp2:e4ent:meas_p95\endcsname{200}
\expandafter\def\csname num:chk:mtp2:e4ent:n\endcsname{92}
\expandafter\def\csname num:chk:mtp2:e4ent:naive_mean\endcsname{145}
\expandafter\def\csname num:chk:mtp2:e4ent:naive_p50\endcsname{135}
\expandafter\def\csname num:chk:mtp2:e4ent:naive_p95\endcsname{210}
\expandafter\def\csname num:chk:mtp2:e4ent:overhead_ms\endcsname{-59}
\expandafter\def\csname num:chk:mtp2:e4ent:pred_agree\endcsname{100.0}
\expandafter\def\csname num:chk:mtp2:e4ent:router_ms\endcsname{0.03}
\expandafter\def\csname num:chk:mtp2:e4ent:sim_acc\endcsname{89.1}
\expandafter\def\csname num:chk:mtp2:e4ent:sim_esc\endcsname{0.0}
\expandafter\def\csname num:chk:mtp2:e4ent:small_agree\endcsname{100.0}
\expandafter\def\csname num:chk:mtp2:e4ent:small_dconf\endcsname{0.00}
\expandafter\def\csname num:chk:mtp2:e4ent:small_frozen_acc\endcsname{89.1}
\expandafter\def\csname num:chk:mtp2:e4ent:small_in_casc_mean\endcsname{139}
\expandafter\def\csname num:chk:mtp2:e4ent:small_live_acc\endcsname{89.1}
\expandafter\def\csname num:chk:mtp2:e4ent:small_mean\endcsname{198}
\expandafter\def\csname num:chk:mtp2:e4ent:small_p50\endcsname{167}
\expandafter\def\csname num:chk:mtp2:e4ent:small_p95\endcsname{538}
\expandafter\def\csname num:chk:mtp2:e4ent:speedup_meas\endcsname{4.4}
\expandafter\def\csname num:chk:mtp2:e4ent:speedup_pred\endcsname{3.1}
\expandafter\def\csname num:chk:mtp2:e4ent:tput_gain\endcsname{3.0}
\expandafter\def\csname num:chk:mtp2:e4ent:x_big8\endcsname{2.65}
\expandafter\def\csname num:chk:mtp2:e4ent:x_bound\endcsname{7.92}
\expandafter\def\csname num:chk:mtp2:e4ent:x_meas\endcsname{7.97}
\expandafter\def\csname num:chk:mtp2:e4ent:x_model\endcsname{8.10}
\expandafter\def\csname num:chk:mtp2:e4ent:x_naive\endcsname{40.49}
\expandafter\def\csname num:chk:mtp2:e4ent:x_small8\endcsname{7.92}
\expandafter\def\csname num:chk:mtp2:e4facts:b8_mean\endcsname{1962}
\expandafter\def\csname num:chk:mtp2:e4facts:big1_agree\endcsname{96.9}
\expandafter\def\csname num:chk:mtp2:e4facts:big1v8_agree\endcsname{97.0}
\expandafter\def\csname num:chk:mtp2:e4facts:big8_agree\endcsname{96.9}
\expandafter\def\csname num:chk:mtp2:e4facts:big_alone_esc_mean\endcsname{503}
\expandafter\def\csname num:chk:mtp2:e4facts:big_frozen_acc\endcsname{88.3}
\expandafter\def\csname num:chk:mtp2:e4facts:big_in_casc_mean\endcsname{467}
\expandafter\def\csname num:chk:mtp2:e4facts:big_live_acc\endcsname{87.8}
\expandafter\def\csname num:chk:mtp2:e4facts:big_mean\endcsname{502}
\expandafter\def\csname num:chk:mtp2:e4facts:big_p50\endcsname{490}
\expandafter\def\csname num:chk:mtp2:e4facts:big_p95\endcsname{555}
\expandafter\def\csname num:chk:mtp2:e4facts:c8_mean\endcsname{909}
\expandafter\def\csname num:chk:mtp2:e4facts:comp_mean\endcsname{293}
\expandafter\def\csname num:chk:mtp2:e4facts:comp_p50\endcsname{96}
\expandafter\def\csname num:chk:mtp2:e4facts:comp_p95\endcsname{630}
\expandafter\def\csname num:chk:mtp2:e4facts:complive_mean\endcsname{293}
\expandafter\def\csname num:chk:mtp2:e4facts:complive_p50\endcsname{96}
\expandafter\def\csname num:chk:mtp2:e4facts:complive_p95\endcsname{630}
\expandafter\def\csname num:chk:mtp2:e4facts:esc_agree\endcsname{100.0}
\expandafter\def\csname num:chk:mtp2:e4facts:live8_acc\endcsname{87.8}
\expandafter\def\csname num:chk:mtp2:e4facts:live8_esc\endcsname{39.9}
\expandafter\def\csname num:chk:mtp2:e4facts:live_acc\endcsname{87.8}
\expandafter\def\csname num:chk:mtp2:e4facts:live_ci\endcsname{[-0.6, +0.6]}
\expandafter\def\csname num:chk:mtp2:e4facts:live_d_big\endcsname{0.0}
\expandafter\def\csname num:chk:mtp2:e4facts:live_esc\endcsname{39.9}
\expandafter\def\csname num:chk:mtp2:e4facts:meas_mean\endcsname{278}
\expandafter\def\csname num:chk:mtp2:e4facts:meas_p50\endcsname{96}
\expandafter\def\csname num:chk:mtp2:e4facts:meas_p95\endcsname{596}
\expandafter\def\csname num:chk:mtp2:e4facts:n\endcsname{1691}
\expandafter\def\csname num:chk:mtp2:e4facts:naive_mean\endcsname{902}
\expandafter\def\csname num:chk:mtp2:e4facts:naive_p50\endcsname{100}
\expandafter\def\csname num:chk:mtp2:e4facts:naive_p95\endcsname{2373}
\expandafter\def\csname num:chk:mtp2:e4facts:overhead_ms\endcsname{-15}
\expandafter\def\csname num:chk:mtp2:e4facts:pred_agree\endcsname{97.6}
\expandafter\def\csname num:chk:mtp2:e4facts:router_ms\endcsname{0.04}
\expandafter\def\csname num:chk:mtp2:e4facts:sim_acc\endcsname{88.2}
\expandafter\def\csname num:chk:mtp2:e4facts:sim_esc\endcsname{39.9}
\expandafter\def\csname num:chk:mtp2:e4facts:small_agree\endcsname{100.0}
\expandafter\def\csname num:chk:mtp2:e4facts:small_dconf\endcsname{0.00}
\expandafter\def\csname num:chk:mtp2:e4facts:small_frozen_acc\endcsname{83.6}
\expandafter\def\csname num:chk:mtp2:e4facts:small_in_casc_mean\endcsname{92}
\expandafter\def\csname num:chk:mtp2:e4facts:small_live_acc\endcsname{83.6}
\expandafter\def\csname num:chk:mtp2:e4facts:small_mean\endcsname{93}
\expandafter\def\csname num:chk:mtp2:e4facts:small_p50\endcsname{91}
\expandafter\def\csname num:chk:mtp2:e4facts:small_p95\endcsname{101}
\expandafter\def\csname num:chk:mtp2:e4facts:speedup_meas\endcsname{1.8}
\expandafter\def\csname num:chk:mtp2:e4facts:speedup_pred\endcsname{1.7}
\expandafter\def\csname num:chk:mtp2:e4facts:tput_gain\endcsname{2.2}
\expandafter\def\csname num:chk:mtp2:e4facts:x_big8\endcsname{4.07}
\expandafter\def\csname num:chk:mtp2:e4facts:x_bound\endcsname{10.22}
\expandafter\def\csname num:chk:mtp2:e4facts:x_meas\endcsname{8.79}
\expandafter\def\csname num:chk:mtp2:e4facts:x_model\endcsname{10.02}
\expandafter\def\csname num:chk:mtp2:e4facts:x_naive\endcsname{27.26}
\expandafter\def\csname num:chk:mtp2:e4facts:x_small8\endcsname{11.67}
\expandafter\def\csname num:sim:e2:clef-flash:acc\endcsname{75.1}
\expandafter\def\csname num:sim:e2:clef-flash:big_acc\endcsname{76.1}
\expandafter\def\csname num:sim:e2:clef-flash:big_ms\endcsname{1077}
\expandafter\def\csname num:sim:e2:clef-flash:ci\endcsname{[-2.0, -0.2]}
\expandafter\def\csname num:sim:e2:clef-flash:d_big\endcsname{-1.0}
\expandafter\def\csname num:sim:e2:clef-flash:esc\endcsname{96}
\expandafter\def\csname num:sim:e2:clef-flash:ms\endcsname{1309}
\expandafter\def\csname num:sim:e2:clef-flash:small_acc\endcsname{30.0}
\expandafter\def\csname num:sim:e2:clef-flash:small_ms\endcsname{273}
\expandafter\def\csname num:sim:e2:clef-flash:speedup\endcsname{0.8}
\expandafter\def\csname num:sim:e2:clef:acc\endcsname{75.1}
\expandafter\def\csname num:sim:e2:clef:big_acc\endcsname{76.1}
\expandafter\def\csname num:sim:e2:clef:big_ms\endcsname{1077}
\expandafter\def\csname num:sim:e2:clef:ci\endcsname{[-2.8, +1.0]}
\expandafter\def\csname num:sim:e2:clef:d_big\endcsname{-1.0}
\expandafter\def\csname num:sim:e2:clef:esc\endcsname{58}
\expandafter\def\csname num:sim:e2:clef:ms\endcsname{1545}
\expandafter\def\csname num:sim:e2:clef:small_acc\endcsname{66.4}
\expandafter\def\csname num:sim:e2:clef:small_ms\endcsname{925}
\expandafter\def\csname num:sim:e2:clef:speedup\endcsname{0.7}
\expandafter\def\csname num:sim:e2:kai:acc\endcsname{75.5}
\expandafter\def\csname num:sim:e2:kai:big_acc\endcsname{76.1}
\expandafter\def\csname num:sim:e2:kai:big_ms\endcsname{1077}
\expandafter\def\csname num:sim:e2:kai:ci\endcsname{[-1.6, +0.2]}
\expandafter\def\csname num:sim:e2:kai:d_big\endcsname{-0.6}
\expandafter\def\csname num:sim:e2:kai:esc\endcsname{97}
\expandafter\def\csname num:sim:e2:kai:ms\endcsname{1111}
\expandafter\def\csname num:sim:e2:kai:small_acc\endcsname{15.3}
\expandafter\def\csname num:sim:e2:kai:small_ms\endcsname{62}
\expandafter\def\csname num:sim:e2:kai:speedup\endcsname{1.0}
\expandafter\def\csname num:sim:e2:lux:acc\endcsname{75.9}
\expandafter\def\csname num:sim:e2:lux:big_acc\endcsname{76.1}
\expandafter\def\csname num:sim:e2:lux:big_ms\endcsname{1077}
\expandafter\def\csname num:sim:e2:lux:ci\endcsname{[-0.6, +0.0]}
\expandafter\def\csname num:sim:e2:lux:d_big\endcsname{-0.2}
\expandafter\def\csname num:sim:e2:lux:esc\endcsname{99}
\expandafter\def\csname num:sim:e2:lux:ms\endcsname{1424}
\expandafter\def\csname num:sim:e2:lux:small_acc\endcsname{31.4}
\expandafter\def\csname num:sim:e2:lux:small_ms\endcsname{357}
\expandafter\def\csname num:sim:e2:lux:speedup\endcsname{0.8}
\expandafter\def\csname num:sim:e2:nox:acc\endcsname{74.8}
\expandafter\def\csname num:sim:e2:nox:big_acc\endcsname{76.1}
\expandafter\def\csname num:sim:e2:nox:big_ms\endcsname{1077}
\expandafter\def\csname num:sim:e2:nox:ci\endcsname{[-2.4, -0.4]}
\expandafter\def\csname num:sim:e2:nox:d_big\endcsname{-1.2}
\expandafter\def\csname num:sim:e2:nox:esc\endcsname{98}
\expandafter\def\csname num:sim:e2:nox:ms\endcsname{1304}
\expandafter\def\csname num:sim:e2:nox:small_acc\endcsname{19.7}
\expandafter\def\csname num:sim:e2:nox:small_ms\endcsname{249}
\expandafter\def\csname num:sim:e2:nox:speedup\endcsname{0.8}
\expandafter\def\csname num:sim:e3:clef-flash:acc\endcsname{84.4}
\expandafter\def\csname num:sim:e3:clef-flash:big_acc\endcsname{84.7}
\expandafter\def\csname num:sim:e3:clef-flash:big_ms\endcsname{528}
\expandafter\def\csname num:sim:e3:clef-flash:ci\endcsname{[-0.9, +0.1]}
\expandafter\def\csname num:sim:e3:clef-flash:d_big\endcsname{-0.3}
\expandafter\def\csname num:sim:e3:clef-flash:esc\endcsname{50}
\expandafter\def\csname num:sim:e3:clef-flash:ms\endcsname{398}
\expandafter\def\csname num:sim:e3:clef-flash:small_acc\endcsname{79.9}
\expandafter\def\csname num:sim:e3:clef-flash:small_ms\endcsname{131}
\expandafter\def\csname num:sim:e3:clef-flash:speedup\endcsname{1.3}
\expandafter\def\csname num:sim:e3:clef:acc\endcsname{84.5}
\expandafter\def\csname num:sim:e3:clef:big_acc\endcsname{84.7}
\expandafter\def\csname num:sim:e3:clef:big_ms\endcsname{528}
\expandafter\def\csname num:sim:e3:clef:ci\endcsname{[-0.9, +0.5]}
\expandafter\def\csname num:sim:e3:clef:d_big\endcsname{-0.2}
\expandafter\def\csname num:sim:e3:clef:esc\endcsname{11}
\expandafter\def\csname num:sim:e3:clef:ms\endcsname{491}
\expandafter\def\csname num:sim:e3:clef:small_acc\endcsname{83.6}
\expandafter\def\csname num:sim:e3:clef:small_ms\endcsname{439}
\expandafter\def\csname num:sim:e3:clef:speedup\endcsname{1.1}
\expandafter\def\csname num:sim:e3:kai:acc\endcsname{85.2}
\expandafter\def\csname num:sim:e3:kai:big_acc\endcsname{84.7}
\expandafter\def\csname num:sim:e3:kai:big_ms\endcsname{528}
\expandafter\def\csname num:sim:e3:kai:ci\endcsname{[+0.0, +1.0]}
\expandafter\def\csname num:sim:e3:kai:d_big\endcsname{+0.5}
\expandafter\def\csname num:sim:e3:kai:esc\endcsname{56}
\expandafter\def\csname num:sim:e3:kai:ms\endcsname{309}
\expandafter\def\csname num:sim:e3:kai:small_acc\endcsname{81.9}
\expandafter\def\csname num:sim:e3:kai:small_ms\endcsname{19}
\expandafter\def\csname num:sim:e3:kai:speedup\endcsname{1.7}
\expandafter\def\csname num:sim:e3:lux:acc\endcsname{84.3}
\expandafter\def\csname num:sim:e3:lux:big_acc\endcsname{84.7}
\expandafter\def\csname num:sim:e3:lux:big_ms\endcsname{528}
\expandafter\def\csname num:sim:e3:lux:ci\endcsname{[-0.9, +0.0]}
\expandafter\def\csname num:sim:e3:lux:d_big\endcsname{-0.5}
\expandafter\def\csname num:sim:e3:lux:esc\endcsname{60}
\expandafter\def\csname num:sim:e3:lux:ms\endcsname{457}
\expandafter\def\csname num:sim:e3:lux:small_acc\endcsname{78.5}
\expandafter\def\csname num:sim:e3:lux:small_ms\endcsname{137}
\expandafter\def\csname num:sim:e3:lux:speedup\endcsname{1.2}
\expandafter\def\csname num:sim:e3:nox:acc\endcsname{84.7}
\expandafter\def\csname num:sim:e3:nox:big_acc\endcsname{84.7}
\expandafter\def\csname num:sim:e3:nox:big_ms\endcsname{528}
\expandafter\def\csname num:sim:e3:nox:ci\endcsname{[-0.3, +0.3]}
\expandafter\def\csname num:sim:e3:nox:d_big\endcsname{+0.0}
\expandafter\def\csname num:sim:e3:nox:esc\endcsname{33}
\expandafter\def\csname num:sim:e3:nox:ms\endcsname{255}
\expandafter\def\csname num:sim:e3:nox:small_acc\endcsname{79.2}
\expandafter\def\csname num:sim:e3:nox:small_ms\endcsname{83}
\expandafter\def\csname num:sim:e3:nox:speedup\endcsname{2.1}
\expandafter\def\csname num:sim:e4ent:clef-flash:acc\endcsname{87.0}
\expandafter\def\csname num:sim:e4ent:clef-flash:big_acc\endcsname{87.0}
\expandafter\def\csname num:sim:e4ent:clef-flash:big_ms\endcsname{616}
\expandafter\def\csname num:sim:e4ent:clef-flash:ci\endcsname{[-5.4, +5.4]}
\expandafter\def\csname num:sim:e4ent:clef-flash:d_big\endcsname{+0.0}
\expandafter\def\csname num:sim:e4ent:clef-flash:esc\endcsname{0}
\expandafter\def\csname num:sim:e4ent:clef-flash:ms\endcsname{193}
\expandafter\def\csname num:sim:e4ent:clef-flash:small_acc\endcsname{87.0}
\expandafter\def\csname num:sim:e4ent:clef-flash:small_ms\endcsname{193}
\expandafter\def\csname num:sim:e4ent:clef-flash:speedup\endcsname{3.2}
\expandafter\def\csname num:sim:e4ent:clef:acc\endcsname{89.1}
\expandafter\def\csname num:sim:e4ent:clef:big_acc\endcsname{87.0}
\expandafter\def\csname num:sim:e4ent:clef:big_ms\endcsname{616}
\expandafter\def\csname num:sim:e4ent:clef:ci\endcsname{[+0.0, +5.4]}
\expandafter\def\csname num:sim:e4ent:clef:d_big\endcsname{+2.2}
\expandafter\def\csname num:sim:e4ent:clef:esc\endcsname{0}
\expandafter\def\csname num:sim:e4ent:clef:ms\endcsname{651}
\expandafter\def\csname num:sim:e4ent:clef:small_acc\endcsname{89.1}
\expandafter\def\csname num:sim:e4ent:clef:small_ms\endcsname{651}
\expandafter\def\csname num:sim:e4ent:clef:speedup\endcsname{0.9}
\expandafter\def\csname num:sim:e4ent:kai:acc\endcsname{84.8}
\expandafter\def\csname num:sim:e4ent:kai:big_acc\endcsname{87.0}
\expandafter\def\csname num:sim:e4ent:kai:big_ms\endcsname{616}
\expandafter\def\csname num:sim:e4ent:kai:ci\endcsname{[-8.7, +4.3]}
\expandafter\def\csname num:sim:e4ent:kai:d_big\endcsname{-2.2}
\expandafter\def\csname num:sim:e4ent:kai:esc\endcsname{7}
\expandafter\def\csname num:sim:e4ent:kai:ms\endcsname{75}
\expandafter\def\csname num:sim:e4ent:kai:small_acc\endcsname{85.9}
\expandafter\def\csname num:sim:e4ent:kai:small_ms\endcsname{37}
\expandafter\def\csname num:sim:e4ent:kai:speedup\endcsname{8.2}
\expandafter\def\csname num:sim:e4ent:lux:acc\endcsname{88.0}
\expandafter\def\csname num:sim:e4ent:lux:big_acc\endcsname{87.0}
\expandafter\def\csname num:sim:e4ent:lux:big_ms\endcsname{616}
\expandafter\def\csname num:sim:e4ent:lux:ci\endcsname{[-2.2, +4.3]}
\expandafter\def\csname num:sim:e4ent:lux:d_big\endcsname{+1.1}
\expandafter\def\csname num:sim:e4ent:lux:esc\endcsname{0}
\expandafter\def\csname num:sim:e4ent:lux:ms\endcsname{213}
\expandafter\def\csname num:sim:e4ent:lux:small_acc\endcsname{88.0}
\expandafter\def\csname num:sim:e4ent:lux:small_ms\endcsname{213}
\expandafter\def\csname num:sim:e4ent:lux:speedup\endcsname{2.9}
\expandafter\def\csname num:sim:e4ent:nox:acc\endcsname{89.1}
\expandafter\def\csname num:sim:e4ent:nox:big_acc\endcsname{87.0}
\expandafter\def\csname num:sim:e4ent:nox:big_ms\endcsname{616}
\expandafter\def\csname num:sim:e4ent:nox:ci\endcsname{[-2.2, +6.5]}
\expandafter\def\csname num:sim:e4ent:nox:d_big\endcsname{+2.2}
\expandafter\def\csname num:sim:e4ent:nox:esc\endcsname{0}
\expandafter\def\csname num:sim:e4ent:nox:ms\endcsname{145}
\expandafter\def\csname num:sim:e4ent:nox:small_acc\endcsname{89.1}
\expandafter\def\csname num:sim:e4ent:nox:small_ms\endcsname{145}
\expandafter\def\csname num:sim:e4ent:nox:speedup\endcsname{4.2}
\expandafter\def\csname num:sim:e4facts:clef-flash:acc\endcsname{87.4}
\expandafter\def\csname num:sim:e4facts:clef-flash:big_acc\endcsname{88.3}
\expandafter\def\csname num:sim:e4facts:clef-flash:big_ms\endcsname{502}
\expandafter\def\csname num:sim:e4facts:clef-flash:ci\endcsname{[-1.4, -0.4]}
\expandafter\def\csname num:sim:e4facts:clef-flash:d_big\endcsname{-0.9}
\expandafter\def\csname num:sim:e4facts:clef-flash:esc\endcsname{66}
\expandafter\def\csname num:sim:e4facts:clef-flash:ms\endcsname{474}
\expandafter\def\csname num:sim:e4facts:clef-flash:small_acc\endcsname{81.6}
\expandafter\def\csname num:sim:e4facts:clef-flash:small_ms\endcsname{143}
\expandafter\def\csname num:sim:e4facts:clef-flash:speedup\endcsname{1.1}
\expandafter\def\csname num:sim:e4facts:clef:acc\endcsname{87.7}
\expandafter\def\csname num:sim:e4facts:clef:big_acc\endcsname{88.3}
\expandafter\def\csname num:sim:e4facts:clef:big_ms\endcsname{502}
\expandafter\def\csname num:sim:e4facts:clef:ci\endcsname{[-1.1, -0.1]}
\expandafter\def\csname num:sim:e4facts:clef:d_big\endcsname{-0.6}
\expandafter\def\csname num:sim:e4facts:clef:esc\endcsname{52}
\expandafter\def\csname num:sim:e4facts:clef:ms\endcsname{741}
\expandafter\def\csname num:sim:e4facts:clef:small_acc\endcsname{76.8}
\expandafter\def\csname num:sim:e4facts:clef:small_ms\endcsname{478}
\expandafter\def\csname num:sim:e4facts:clef:speedup\endcsname{0.7}
\expandafter\def\csname num:sim:e4facts:kai:acc\endcsname{87.8}
\expandafter\def\csname num:sim:e4facts:kai:big_acc\endcsname{88.3}
\expandafter\def\csname num:sim:e4facts:kai:big_ms\endcsname{502}
\expandafter\def\csname num:sim:e4facts:kai:ci\endcsname{[-1.0, -0.1]}
\expandafter\def\csname num:sim:e4facts:kai:d_big\endcsname{-0.5}
\expandafter\def\csname num:sim:e4facts:kai:esc\endcsname{91}
\expandafter\def\csname num:sim:e4facts:kai:ms\endcsname{478}
\expandafter\def\csname num:sim:e4facts:kai:small_acc\endcsname{63.6}
\expandafter\def\csname num:sim:e4facts:kai:small_ms\endcsname{21}
\expandafter\def\csname num:sim:e4facts:kai:speedup\endcsname{1.0}
\expandafter\def\csname num:sim:e4facts:lux:acc\endcsname{88.1}
\expandafter\def\csname num:sim:e4facts:lux:big_acc\endcsname{88.3}
\expandafter\def\csname num:sim:e4facts:lux:big_ms\endcsname{502}
\expandafter\def\csname num:sim:e4facts:lux:ci\endcsname{[-0.7, +0.2]}
\expandafter\def\csname num:sim:e4facts:lux:d_big\endcsname{-0.2}
\expandafter\def\csname num:sim:e4facts:lux:esc\endcsname{49}
\expandafter\def\csname num:sim:e4facts:lux:ms\endcsname{387}
\expandafter\def\csname num:sim:e4facts:lux:small_acc\endcsname{86.3}
\expandafter\def\csname num:sim:e4facts:lux:small_ms\endcsname{142}
\expandafter\def\csname num:sim:e4facts:lux:speedup\endcsname{1.3}
\expandafter\def\csname num:sim:e4facts:nox:acc\endcsname{88.2}
\expandafter\def\csname num:sim:e4facts:nox:big_acc\endcsname{88.3}
\expandafter\def\csname num:sim:e4facts:nox:big_ms\endcsname{502}
\expandafter\def\csname num:sim:e4facts:nox:ci\endcsname{[-0.7, +0.5]}
\expandafter\def\csname num:sim:e4facts:nox:d_big\endcsname{-0.1}
\expandafter\def\csname num:sim:e4facts:nox:esc\endcsname{40}
\expandafter\def\csname num:sim:e4facts:nox:ms\endcsname{295}
\expandafter\def\csname num:sim:e4facts:nox:small_acc\endcsname{83.6}
\expandafter\def\csname num:sim:e4facts:nox:small_ms\endcsname{94}
\expandafter\def\csname num:sim:e4facts:nox:speedup\endcsname{1.7}
\expandafter\def\csname num:sim:evals\endcsname{628{,}880}
\expandafter\def\csname num:sim:seconds\endcsname{17.7}
\expandafter\def\csname num:xcfg:e2:mtp0:mtp0apc0\endcsname{92.2}
\expandafter\def\csname num:xcfg:e2:mtp2:mtp0\endcsname{95.0}
\expandafter\def\csname num:xcfg:e2:mtp2:mtp0apc0\endcsname{90.7}
\expandafter\def\csname num:xcfg:e3:mtp0:mtp0apc0\endcsname{99.2}
\expandafter\def\csname num:xcfg:e3:mtp2:mtp0\endcsname{99.5}
\expandafter\def\csname num:xcfg:e3:mtp2:mtp0apc0\endcsname{99.2}
\expandafter\def\csname num:xcfg:e4ent:mtp0:mtp0apc0\endcsname{98.9}
\expandafter\def\csname num:xcfg:e4ent:mtp2:mtp0\endcsname{98.9}
\expandafter\def\csname num:xcfg:e4ent:mtp2:mtp0apc0\endcsname{97.8}
\expandafter\def\csname num:xcfg:e4facts:mtp0:mtp0apc0\endcsname{97.2}
\expandafter\def\csname num:xcfg:e4facts:mtp2:mtp0\endcsname{98.5}
\expandafter\def\csname num:xcfg:e4facts:mtp2:mtp0apc0\endcsname{97.2}

\newcommand{\N}[1]{\csname num:#1\endcsname}
\newcommand{\TODO}[1]{\textbf{[#1]}}

\title{Simulate Before You Serve:\\
Predicting a Decision-Model Cascade from Frozen Outputs,\\
and Checking the Prediction on Real Hardware}

\author{Kostadis Roussos\\
Independent researcher\\
\texttt{kostadis@gmail.com}}

\date{October 2026}

\begin{document}
\maketitle

\begin{abstract}
A cascade answers with a cheap model and escalates uncertain items to an expensive one. Choosing the
cheap model, the gate and the threshold usually means deploying and measuring each candidate. We show
that when every candidate has already been run on a frozen evaluation set, the cascade can instead be
\emph{replayed} from the saved per-item answers, confidences and request times, and we check that
replay against the real system. The setting is a working pipeline for tabletop role-playing campaign
records. Five non-generative decision models are candidate first stages in front of a generative model
(\texttt{qwen3.8-flash-next}) on four judgment tasks labelled by a game master, served on two NVIDIA
DGX Spark (GB10) systems, with the goal of the generative model's accuracy at lower latency. The
replay evaluated the whole design space in \N{sim:seconds}~s of CPU time. It chose Decision-2.0-Nox-4B
as the first stage for two tasks, Nox alone for a third, and no cascade for the fourth. We then ran that
cascade live with thresholds frozen in advance. It held the generative model's accuracy (the
cascade-minus-qwen difference was predicted within 0.1 points on the three tasks with escalation)
while answering
\N{chk:mtp2:e3:speedup_meas}$\times$ and \N{chk:mtp2:e4facts:speedup_meas}$\times$ faster
single-stream and with \N{chk:mtp2:e3:tput_gain}$\times$ and \N{chk:mtp2:e4facts:tput_gain}$\times$
the throughput under 8-way load. Escalation decisions were predicted exactly, because the decision
model is bit-for-bit deterministic; the generative model at temperature~0 was not, changing 1--7\% of
its answers with load and serving configuration. Latency predictions were 3--4$\times$ wrong when the
replay mixed load regimes and within 5--9\% when it did not. A two-station queueing model predicted
throughput within 3\% when one stage dominated and 14--25\% high when the stages were balanced. A
serving ablation found that, with speculative decoding off, prefix-cache hits on this model's serving
image were 38\% slower and changed answers, a hidden variable for any replay.
\end{abstract}

\section{Introduction}

A practitioner who wants to put a cheap model in front of an expensive one faces a design space
before any deployment question: which cheap model, which gate, which threshold, for which task.
Each point in that space can be measured by deploying it, at the cost of GPU time, a serving
change and a benchmark run per point. Many such points, however, are already determined by
experiments that have been run. If every candidate model has been run on every item of a frozen
evaluation set, and inference is deterministic, then the answer, the confidence and the request time
of any cascade on any item are already recorded. The cascade can be \emph{replayed} rather than run.

This paper tests that idea on a working pipeline. A companion study~\cite{companion} ran six
non-generative decision models and a generative model, \texttt{qwen3.8-flash-next}, on a frozen
dataset of one game master's rulings over tabletop role-playing transcripts. Its practical verdict was
that the generative model was the most accurate on most judgment steps and the decision models were
an order of magnitude faster. vLLM Semantic Router's \emph{System One Auto}~\cite{sysoneauto}
suggests combining the two, by answering with a small decision model and escalating uncertain items.
Our goal is narrow: \emph{the generative model's accuracy at lower latency}. We study it twice:
\begin{enumerate}
  \item \textbf{Analytically}, by replaying the frozen outputs: five candidate first stages, four
        tasks, a held-out threshold that targets parity with the generative model, and latency and
        throughput composed from measured service times. The whole design space evaluates in seconds
        of CPU time.
  \item \textbf{Empirically}, by running the cascade the replay recommends, Decision-2.0-Nox-4B in
        front of qwen, as a real two-box router on two NVIDIA DGX Spark (GB10) systems, with every
        threshold frozen before the first live request, single-stream and under 8-way load. We then
        repeat the qwen arms under three serving configurations, to find out what breaks the
        replay's central assumption.
\end{enumerate}

\paragraph{Findings.} (i) The replay picked the right design. Nox-4B $\rightarrow$ qwen held qwen's
accuracy on triage and fact typing while answering 1.8--2.2$\times$ faster single-stream and
2.2--2.4$\times$ faster under load; Nox alone sufficed for entity typing; spelling should not be
cascaded. The other four first stages were each too weak or too slow. (ii) Escalation decisions were
predicted \emph{exactly}, because the decision model is bit-for-bit deterministic. (iii) The
cascade-minus-qwen accuracy difference was predicted within 0.1 points, although the generative
model at temperature~0 did not reproduce 1--7\% of its own answers. (iv) Latency predictions were
3--4$\times$ wrong when the replay mixed load regimes and within 5--9\% when it did not. A
two-station queueing model predicted 8-client throughput within 3\% when one stage dominates and
overpredicted it by 14--25\% when the two stages are balanced. (v) Turning speculative decoding off
did not make the generative model reproducible; its answers depend on load and serving configuration
together. With speculative decoding off and prefix caching on, cache hits on this model's serving image
became 38\% slower and changed answers, which made the replay underpredict cascade latency by a
quarter.

\paragraph{Why it matters.} The replay turned the design question into a table computed in under a
minute; the live check took several GPU-hours across two machines and a serving restart for each
configuration tested. Where a replay is trustworthy, and where it quietly is not, is the practical
content of this paper.

\section{Background}

\subsection{The companion study and its frozen dataset}
A companion paper~\cite{companion} evaluated six non-generative \emph{decision models} against a
generative model on judgment steps of a pipeline that turns recorded tabletop role-playing sessions
into campaign records. A decision model reads a \texttt{state} once and a small head returns a
probability for every option of every typed question (\texttt{choice}, \texttt{noul},
\texttt{score}) through the SystemOne interface~\cite{typesafe}; no text is generated. Labels are one
game master's (GM's) own past rulings, rendered once into a read-only, versioned dataset with every
model-facing \texttt{state} already built, so every model was run against the same bytes.
That study treated the models as alternatives. Its conclusion for the choice tasks was mixed: the
generative model, \texttt{qwen3.8-flash-next}, was usually the most accurate, while the decision
models were 10--25$\times$ faster and sometimes as accurate. This paper asks whether the two can be
\emph{combined}.

We reuse four of its tasks, unchanged:
\begin{description}
  \item[E2 spelling] (497 items) choose the canonical spelling of a speech-recognition-mangled name
        from a candidate list, or \texttt{LEAVE}/\texttt{OTHER}.
  \item[E3 triage] (864 tokens) is an unknown capitalised token a campaign name (forward it to the
        spelling step) or not (filler, rules term, real-world chatter)? 84 are real names. Here we
        score it as binary accuracy and separately count \emph{missed names}, the error that
        matters: a dropped name is a silent deletion nothing downstream can recover.
  \item[E4 fact typing] (1{,}691 facts) file an extracted fact under one of eight types.
  \item[E4 entity typing] (92 GM merge rulings) the entity type of a dossier subject, five types.
\end{description}

\subsection{Cascades and System One Auto}
A cascade answers with a cheap model and re-asks an expensive one only when the cheap model is
unsure~\cite{frugalgpt,routellm}. vLLM Semantic Router's \emph{System One Auto}~\cite{sysoneauto}
applies the idea to decision models: Decision-2.0-Kai-0.6B answers first, and a gate on its top
probability ($\max(p,1-p)$ for yes/no questions) escalates the request to Decision-2.0-Vega-27B.
On 231 public benchmark items, with a threshold frozen on a separate calibration set, the cascade
scored 82.3\% between Kai (65.8\%) and Vega (87.9\%) while escalating 46\% of requests. Two
differences matter for us. Our large model is \emph{generative} and reached through a chat endpoint,
so the two stages do not share an interface or a runtime; and our goal is narrower: \emph{the large
model's accuracy at lower latency}, not a point between the two.

\section{Method}

\subsection{The cascade}
For each item the small model $S$ (a decision model) returns an answer $s$ and a gate confidence
$c$: the top option probability for \texttt{choice} questions, and $\max(p, 1-p)$ for the binary
E3 reading, where $p = P(\texttt{campaign})$. If $c < t$ the same question goes to the large model
$B$, prompted exactly as in the companion study (the option list rendered into a chat prompt,
JSON answer parsed), and the cascade answers with $B$'s answer; otherwise it answers $s$. $B$ is
\texttt{qwen3.8-flash-next} with reasoning disabled, on one DGX Spark (spark1); every $S$ runs on
the second (spark2). The two boxes are on the same LAN as the client.

\subsection{The replay model}
\label{sec:replay}
The companion study already ran every model on every item, at temperature~0 (generative) or with a
deterministic head (decision models), and kept the per-item answers, confidences and request
times. If inference is deterministic, the cascade's answer on item $i$ at threshold $t$ is a pure
function of those saved records:
\[
  \hat a_i(t) = \begin{cases} b_i & c_i < t \\ s_i & \text{otherwise}\end{cases},\qquad
  \hat T_i(t) = T^S_i + [c_i < t]\,T^B_i ,
\]
so accuracy, escalation rate and the full latency distribution of any (small model, threshold)
pair can be computed without running a model. We call this the \emph{replay model}. It makes two
assumptions that the live experiment tests: \textbf{(A1) determinism}, that a model asked again
returns the saved answer (and, for $S$, the saved confidence, since the gate depends on it); and
\textbf{(A2) composable service times}, that a cascade request costs one $S$ request plus, if
escalated, one $B$ request, each as long as when measured alone in the same load regime.
A2 is a statement about load: the saved $T^B$ came from an 8-way concurrent run and $T^S$ from
single-stream runs, so a naive replay mixes regimes. We therefore report latency from the replay
with $T^B$ re-measured single-stream (Section~\ref{sec:live}), and keep the naive figure only to
show the size of that error.

For concurrent load the composition is no longer additive. We model the $c$-client closed loop as
two stations visited with probabilities $1$ ($S$) and $p$ ($B$, the escalation rate), each with a
response time $R(n)$ that depends on the number $n$ of requests in it, linearly interpolated
between its measured single-stream and 8-way means. Little's law~\cite{little1961} at each station,
$n_S = X R_S(n_S)$ and $n_B = X p R_B(n_B)$ with $n_S + n_B = c$, is solved for the throughput $X$,
in the spirit of mean-value analysis with load-dependent stations~\cite{reiser1980}. We compare it
with two simpler predictions: no contention, $X = c / \mathbb{E}[\hat T]$, and the bottleneck bound
$X \le \min(X_S^{\max}, X_B^{\max}/p)$.

\subsection{Choosing the threshold}
\label{sec:gate}
Our objective is parity with $B$ at the lowest cost. On a training set the \emph{parity gate} is the
lowest threshold whose cascade accuracy is at least $B$'s accuracy alone (and, on E3, that misses
no more names than $B$); because escalating everything reproduces $B$ exactly, such a threshold
always exists. Thresholds are fitted 2-fold by group: sessions (E2, E3), chapters (E4 facts) or
subjects (E4 entities) are sorted and split alternately, the gate is fitted on one half and applied
to the other, so no item's own label influences its threshold. Allowing a tolerance $\delta>0$
below parity did not generalise in a pilot: a $\delta = 0.5$ point budget on the training half
became a 1.3--1.6 point loss on the test half for E4 facts, so we use $\delta = 0$ throughout.

\subsection{The live experiment}
\label{sec:live}
For the pair Nox-4B $\rightarrow$ qwen we froze the held-out parity gates for every item before any
live run (file SHA-256 \texttt{aa21ed3f\ldots}), together with the replay's predicted escalation
and answer per item. A small router (\texttt{cascade\_live.py}) then ran, against the real endpoints
and using the same question definitions, prompts and parsers as the companion study:
$S$ alone, $B$ alone and the cascade, each single-stream ($c = 1$) and with 8 concurrent clients
($c = 8$), on all four tasks. Every request's wall time is measured at the client, so cascade
latency includes the router and both network hops. Live and replay results are compared item by
item.

\section{Results from the replay model}
\label{sec:sim}

Table~\ref{tab:design} is the replay model's answer to the design question: for each of five
decision models in front of qwen, and each task, the held-out parity gate's accuracy, escalation
rate and mean latency. Accuracy comes from the frozen per-item answers; latency composes each
decision model's single-stream request times from the companion study with qwen's request times
re-measured single-stream on the same items (Section~\ref{sec:live}). Producing the whole table,
plus the threshold sweeps behind Figure~\ref{fig:pareto}, took \N{sim:seconds}~s of CPU time and
\N{sim:evals} cascade-item evaluations, and no GPU time at all.

\begin{table}[t]
\centering
\small
\begin{tabular}{llrrrrrrr}
\toprule
Task & Small & \multicolumn{3}{c}{Accuracy (\%)} & $\Delta$ vs qwen (pt) & Esc. & Mean & Speed-up \\
\cmidrule(lr){3-5}
 & model & small & qwen & cascade & [95\% CI] & (\%) & (ms) & vs qwen \\
\midrule
\multirow{5}{*}{E3 triage} & nox & 79.2 & 84.7 & 84.7 (6/6) & +0.0 [-0.3, +0.3] & 33 & 255 & 2.1$\times$ \\
 & kai & 81.9 & 84.7 & 85.2 (7/6) & +0.5 [+0.0, +1.0] & 56 & 309 & 1.7$\times$ \\
 & clef-flash & 79.9 & 84.7 & 84.4 (6/6) & -0.3 [-0.9, +0.1] & 50 & 398 & 1.3$\times$ \\
 & clef & 83.6 & 84.7 & 84.5 (6/6) & -0.2 [-0.9, +0.5] & 11 & 491 & 1.1$\times$ \\
 & lux & 78.5 & 84.7 & 84.3 (6/6) & -0.5 [-0.9, +0.0] & 60 & 457 & 1.2$\times$ \\
\midrule
\multirow{5}{*}{E4 fact typing} & nox & 83.6 & 88.3 & 88.2 & -0.1 [-0.7, +0.5] & 40 & 295 & 1.7$\times$ \\
 & kai & 63.6 & 88.3 & 87.8 & -0.5 [-1.0, -0.1] & 91 & 478 & 1.0$\times$ \\
 & clef-flash & 81.6 & 88.3 & 87.4 & -0.9 [-1.4, -0.4] & 66 & 474 & 1.1$\times$ \\
 & clef & 76.8 & 88.3 & 87.7 & -0.6 [-1.1, -0.1] & 52 & 741 & 0.7$\times$ \\
 & lux & 86.3 & 88.3 & 88.1 & -0.2 [-0.7, +0.2] & 49 & 387 & 1.3$\times$ \\
\midrule
\multirow{5}{*}{E4 entity typing} & nox & 89.1 & 87.0 & 89.1 & +2.2 [-2.2, +6.5] & 0 & 145 & 4.2$\times$ \\
 & kai & 85.9 & 87.0 & 84.8 & -2.2 [-8.7, +4.3] & 7 & 75 & 8.2$\times$ \\
 & clef-flash & 87.0 & 87.0 & 87.0 & +0.0 [-5.4, +5.4] & 0 & 193 & 3.2$\times$ \\
 & clef & 89.1 & 87.0 & 89.1 & +2.2 [+0.0, +5.4] & 0 & 651 & 0.9$\times$ \\
 & lux & 88.0 & 87.0 & 88.0 & +1.1 [-2.2, +4.3] & 0 & 213 & 2.9$\times$ \\
\midrule
\multirow{5}{*}{E2 spelling} & nox & 19.7 & 76.1 & 74.8 & -1.2 [-2.4, -0.4] & 98 & 1304 & 0.8$\times$ \\
 & kai & 15.3 & 76.1 & 75.5 & -0.6 [-1.6, +0.2] & 97 & 1111 & 1.0$\times$ \\
 & clef-flash & 30.0 & 76.1 & 75.1 & -1.0 [-2.0, -0.2] & 96 & 1309 & 0.8$\times$ \\
 & clef & 66.4 & 76.1 & 75.1 & -1.0 [-2.8, +1.0] & 58 & 1545 & 0.7$\times$ \\
 & lux & 31.4 & 76.1 & 75.9 & -0.2 [-0.6, +0.0] & 99 & 1424 & 0.8$\times$ \\
\bottomrule
\end{tabular}

\caption{Replay-model predictions for every decision model in front of qwen (reasoning off), with the
held-out parity gate. ``Esc.''\ is the share of items sent to qwen. Mean latency is per item,
single-stream. For E3 the cascade column also gives missed real names (cascade/qwen) out of 84.
CIs are paired bootstrap over items. The speed-up is qwen's mean single-stream latency on the same
items divided by the cascade's.}
\label{tab:design}
\end{table}

\paragraph{Nox-4B is the only small model that pays off broadly.} It holds parity on E3
(\N{sim:e3:nox:acc}\% against qwen's \N{sim:e3:nox:big_acc}\%, the same 6 missed names) with
\N{sim:e3:nox:esc}\% escalated, a predicted \N{sim:e3:nox:speedup}$\times$ speed-up; and on E4 fact
typing (\N{sim:e4facts:nox:acc}\% against \N{sim:e4facts:nox:big_acc}\%, CI \N{sim:e4facts:nox:ci})
with \N{sim:e4facts:nox:esc}\% escalated, \N{sim:e4facts:nox:speedup}$\times$. On E4 entity typing it
needs no escalation at all: alone it is at least as accurate as qwen (\N{sim:e4ent:nox:acc}\% vs
\N{sim:e4ent:nox:big_acc}\%), so the ``cascade'' is Nox alone at \N{sim:e4ent:nox:speedup}$\times$.

\paragraph{The other small models fail for two different reasons.} Kai-0.6B is the fastest model
(\N{sim:e4facts:kai:small_ms}~ms) but too weak on fact typing (\N{sim:e4facts:kai:small_acc}\%): the
gate must send \N{sim:e4facts:kai:esc}\% of items to qwen, and the cascade is no faster than qwen.
Clef-27B is the opposite: accurate enough to escalate rarely on E3 (\N{sim:e3:clef:esc}\%), but its
own request (\N{sim:e3:clef:small_ms}~ms) costs most of a qwen call (\N{sim:e3:clef:big_ms}~ms), so it
saves almost nothing; on fact typing the Clef cascade is \emph{slower} than qwen alone. A first
stage only helps if it is both much cheaper than the second and confident on a large share of
items; the replay makes that trade-off visible per task, before anything is deployed.

\paragraph{Spelling (E2) should not be cascaded.} Every decision model is far below qwen on choosing
a canonical spelling (15--66\% against \N{sim:e2:nox:big_acc}\%), so the parity gate escalates
58--99\% of items and every cascade is slower than qwen alone. This reproduces the companion
study's finding that string-level judgments are where decision models fail.

\paragraph{No cascade beats qwen.} Picking the threshold on the test items themselves (an upper
bound) raises cascade accuracy by at most a few tenths of a point over qwen on E3 and E4 facts.
The ceiling is set by the items both models get wrong: 124 of 864 on E3 and 140 of 1{,}691 on E4
facts. A cascade picks between
two answers; it cannot produce a right answer neither model has. The binary E3 reading has one trap,
noted here because the replay exposed it: a gate fitted for plain accuracy (rather than parity)
finds a Kai threshold that ``beats'' qwen by keeping confident not-a-name answers while missing
about half the real names. On this task that is the expensive error, which is why the parity gate
also constrains missed names.

\begin{figure}[t]
\centering
\begin{tikzpicture}
\begin{groupplot}[group style={group size=2 by 1, horizontal sep=1.6cm},
    width=0.5\textwidth, height=0.42\textwidth, xlabel={mean latency per item (ms)},
    ylabel={accuracy (\%)}, grid=major, grid style={gray!20}, legend style={font=\scriptsize},
    tick label style={font=\scriptsize}, label style={font=\small}, title style={font=\small}]
\nextgroupplot[title={E3 triage}, xmin=0, xmax=900, ymin=76, ymax=88, legend to name=paretoleg, legend columns=6]
\addplot[thick, blue] table[x=ms, y=acc] {gen/sweep-e3-nox.dat}; \addlegendentry{Nox-4B}
\addplot[thick, orange] table[x=ms, y=acc] {gen/sweep-e3-kai.dat}; \addlegendentry{Kai-0.6B}
\addplot[thick, teal] table[x=ms, y=acc] {gen/sweep-e3-clef-flash.dat}; \addlegendentry{clef-flash}
\addplot[thick, red] table[x=ms, y=acc] {gen/sweep-e3-clef.dat}; \addlegendentry{Clef-27B}
\addplot[thick, violet] table[x=ms, y=acc] {gen/sweep-e3-lux.dat}; \addlegendentry{Lux-9B}
\addplot[only marks, mark=*, black] coordinates {(\N{sim:e3:nox:big_ms}, \N{sim:e3:nox:big_acc})}; \addlegendentry{qwen alone}
\nextgroupplot[title={E4 fact typing}, xmin=0, xmax=900, ymin=60, ymax=90]
\addplot[thick, blue] table[x=ms, y=acc] {gen/sweep-e4facts-nox.dat};
\addplot[thick, orange] table[x=ms, y=acc] {gen/sweep-e4facts-kai.dat};
\addplot[thick, teal] table[x=ms, y=acc] {gen/sweep-e4facts-clef-flash.dat};
\addplot[thick, red] table[x=ms, y=acc] {gen/sweep-e4facts-clef.dat};
\addplot[thick, violet] table[x=ms, y=acc] {gen/sweep-e4facts-lux.dat};
\addplot[only marks, mark=*, black] coordinates {(\N{sim:e4facts:nox:big_ms}, \N{sim:e4facts:nox:big_acc})};
\end{groupplot}
\end{tikzpicture}
\\[2pt]\pgfplotslegendfromname{paretoleg}
\caption{Replay-model accuracy--latency curves, sweeping the gate threshold from 0 (small model
alone, left end of each curve) to 1 (everything escalated, converging on the qwen point). All items,
not held out; single-stream latency. A useful cascade is a curve that reaches qwen's accuracy well
to the left of the qwen point. Kai's E3 peak above qwen is the plain-accuracy artifact discussed in
the text: it keeps confident not-a-name answers and misses real names.}
\label{fig:pareto}
\end{figure}

\section{Results from the live cascade}
\label{sec:liveres}

Table~\ref{tab:live} reports the Nox-4B $\rightarrow$ qwen cascade run for real with the
pre-registered gates, against qwen alone run live on the same items in the same session, with qwen
in its deployed configuration (MTP-2 speculative decoding, prefix caching on).

\begin{table}[t]
\centering
\small
\setlength{\tabcolsep}{4pt}
\resizebox{\textwidth}{!}{\begin{tabular}{lrrrrrrrrrrr}
\toprule
 & \multicolumn{4}{c}{Accuracy (\%), $c=1$} & Esc. & \multicolumn{3}{c}{Mean latency, $c=1$} & \multicolumn{3}{c}{Throughput, $c=8$ (items/s)} \\
\cmidrule(lr){2-5}\cmidrule(lr){7-9}\cmidrule(lr){10-12}
Task & Nox & qwen & cascade & $\Delta$ [95\% CI] & (\%) & qwen & cascade & speed-up & qwen & cascade & gain \\
\midrule
E3 triage & 79.2 & 84.4 & 84.5 & +0.1 [+0.0, +0.3] & 33 & 528 & 241 & 2.19$\times$ & 4.08 & 9.84 & 2.41$\times$ \\
E4 fact typing & 83.6 & 87.8 & 87.8 & +0.0 [-0.6, +0.6] & 40 & 502 & 278 & 1.80$\times$ & 4.07 & 8.79 & 2.16$\times$ \\
E4 entity typing & 89.1 & 88.0 & 89.1 & +1.1 [-2.2, +5.4] & 0 & 616 & 139 & 4.43$\times$ & 2.65 & 7.97 & 3.00$\times$ \\
E2 spelling & 19.7 & 75.9 & 74.6 & -1.2 [-2.4, -0.4] & 98 & 1077 & 1227 & 0.88$\times$ & 1.74 & 1.72 & 0.99$\times$ \\
\bottomrule
\end{tabular}
}
\caption{The live cascade (Nox-4B on spark2 $\rightarrow$ qwen on spark1) against qwen alone,
measured in the same session. Accuracy and mean latency per item are single-stream ($c=1$). $\Delta$
is cascade minus qwen with a paired bootstrap CI. Throughput uses 8 concurrent clients over the whole
task. Latency is wall time at the client, including both network hops.}
\label{tab:live}
\end{table}

\paragraph{Accuracy parity held.} On E3 the cascade matched qwen (\N{chk:mtp2:e3:live_acc}\% vs
\N{chk:mtp2:e3:big_live_acc}\%) and missed the same number of real names
(\N{chk:mtp2:e3:live_missed} vs \N{chk:mtp2:e3:big_live_missed}). On E4 fact typing it was
exactly level (\N{chk:mtp2:e4facts:live_acc}\% vs \N{chk:mtp2:e4facts:big_live_acc}\%, CI
\N{chk:mtp2:e4facts:live_ci}). On E4 entity typing Nox alone answered everything and was at least as
accurate as qwen. On E2 the cascade trailed qwen ($\Delta = \N{chk:mtp2:e2:live_d_big}$~points), a
loss the replay had predicted to the tenth of a point (next section).

\paragraph{Latency and throughput improved where the replay said they would.} Single-stream, the
cascade answered E3 \N{chk:mtp2:e3:speedup_meas}$\times$ and E4 facts
\N{chk:mtp2:e4facts:speedup_meas}$\times$ faster than qwen on average
(\N{chk:mtp2:e3:meas_mean} vs \N{chk:mtp2:e3:big_mean}~ms and \N{chk:mtp2:e4facts:meas_mean} vs
\N{chk:mtp2:e4facts:big_mean}~ms), and E4 entities \N{chk:mtp2:e4ent:speedup_meas}$\times$. The
distribution is bimodal: the median request is a Nox-only answer
(\N{chk:mtp2:e3:meas_p50}~ms on E3), while the 95th percentile is an escalated one
(\N{chk:mtp2:e3:meas_p95}~ms, against \N{chk:mtp2:e3:big_p95}~ms for qwen alone). An escalated
request costs both stages, so the cascade's tail is slightly \emph{worse} than qwen's, as
System One Auto also reported~\cite{sysoneauto}. Under 8-way load the gain is larger:
\N{chk:mtp2:e3:tput_gain}$\times$ (E3), \N{chk:mtp2:e4facts:tput_gain}$\times$ (E4 facts) and
\N{chk:mtp2:e4ent:tput_gain}$\times$ (E4 entities) the throughput of qwen alone, because escalating
33--40\% of items removes most of qwen's queue. On E2, where 98\% of items escalate, the cascade
was slower than qwen alone single-stream (\N{chk:mtp2:e2:speedup_meas}$\times$) and no faster under
load (\N{chk:mtp2:e2:tput_gain}$\times$), as the replay predicted.

\section{Prediction versus measurement}
\label{sec:check}

\subsection{Accuracy: exact where the models are deterministic}
Table~\ref{tab:check} compares the replay's item-level predictions with the live run.

\begin{table}[t]
\centering
\small
\setlength{\tabcolsep}{4pt}
\resizebox{\textwidth}{!}{\begin{tabular}{lrrrrrrrrrr}
\toprule
 & \multicolumn{2}{c}{Escalation} & \multicolumn{2}{c}{Answers} & \multicolumn{2}{c}{$\Delta$ acc.\ vs qwen (pt)} & \multicolumn{4}{c}{Mean latency, $c=1$ (ms)} \\
\cmidrule(lr){2-3}\cmidrule(lr){4-5}\cmidrule(lr){6-7}\cmidrule(lr){8-11}
Task & pred. & agree & agree & qwen repro & pred. & meas. & meas. & naive & replay & error \\
\midrule
E3 triage & 33.1\% & 100\% & 99.3\% & 99.0\% & +0.0 & +0.1 & 241 & 746 & 263 & +9\% \\
E4 fact typing & 39.9\% & 100\% & 97.6\% & 96.9\% & -0.1 & +0.0 & 278 & 902 & 293 & +5\% \\
E4 entity typing & 0.0\% & 100\% & 100.0\% & 98.9\% & +2.2 & +1.1 & 139 & 145 & 198 & +42\% \\
E2 spelling & 98.0\% & 100\% & 93.2\% & 93.2\% & -1.2 & -1.2 & 1227 & 5080 & 1336 & +9\% \\
\bottomrule
\end{tabular}
}
\caption{Replay prediction against the live cascade ($c=1$, qwen deployed configuration).
\emph{Escalation}: predicted share escalated, and the share of items where the live router made
the predicted escalation decision. \emph{Answers}: share of items where the live cascade's answer
equals the predicted one, and share where live qwen alone reproduces its own frozen answer.
\emph{$\Delta$ acc.}: cascade minus qwen, predicted (frozen answers) and measured (live run, same
session). \emph{Latency}: measured mean; the naive replay using the frozen request times (qwen's were
8-way concurrent); the replay with single-stream service times measured live; and its relative error.}
\label{tab:check}
\end{table}

\paragraph{Assumption A1 holds for the decision model, not for the generative one.} Nox
reproduced every frozen answer and every frozen confidence bit for bit (maximum absolute difference exactly zero across all \N{chk:mtp2:e4facts:n}+\N{chk:mtp2:e3:n}+\N{chk:mtp2:e4ent:n}+\N{chk:mtp2:e2:n}
items), so the live router escalated \emph{exactly} the predicted items on every task. qwen, at
temperature~0 in its deployed configuration, did not: it reproduced its own frozen answer on
\N{chk:mtp2:e3:big1_agree}\% of E3 items, \N{chk:mtp2:e4facts:big1_agree}\% of E4 facts and only
\N{chk:mtp2:e2:big1_agree}\% of E2, where answers are longer. Every live-versus-predicted answer
mismatch in the cascade, in every serving configuration we ran, is an escalated item on which qwen
changed its answer; none comes from the gate. On E4 entity typing, where nothing escalates, the
predicted and measured differences from qwen (+2.2 and +1.1 points) differ only because qwen alone
changed one of its 92 answers. qwen's single-stream and 8-way answers
from the same session also agree on only \N{chk:mtp2:e4facts:big1v8_agree}\% of E4 facts. Section~\ref{sec:ablation}
examines where it comes from.

\paragraph{The replay predicts the cascade's \emph{effect} better than its level.} qwen's drift moves the
absolute accuracy of both qwen and the cascade together, so the replay's absolute figure for E4
facts (\N{chk:mtp2:e4facts:sim_acc}\%) is half a point above the live one
(\N{chk:mtp2:e4facts:live_acc}\%). But the quantity the design decision rests on, the difference
from qwen alone, was predicted within 0.1 points on E3, E4 facts and E2, including E2's predicted
\emph{loss}. Both arms inherit the same qwen, so its noise largely cancels in the comparison.

\subsection{Latency: right only in the right load regime}
The naive replay, which uses the request times recorded in the companion study, predicted mean
cascade latencies 3--4$\times$ too high (\N{chk:mtp2:e3:naive_mean} against
\N{chk:mtp2:e3:meas_mean}~ms on E3), because qwen's recorded times came from an 8-way concurrent run
while the cascade was measured single-stream. With service times measured in the target regime the
same replay is within 5--9\% on the three tasks with escalation, and predicts the speed-ups
(\N{chk:mtp2:e3:speedup_pred}$\times$ and \N{chk:mtp2:e4facts:speedup_pred}$\times$ predicted,
\N{chk:mtp2:e3:speedup_meas}$\times$ and \N{chk:mtp2:e4facts:speedup_meas}$\times$ measured).
The residual error is not router overhead, which measured \N{chk:mtp2:e3:router_ms}~ms per request;
it is drift in the models' own service times between runs. qwen answered the escalated E4-facts
items in \N{chk:mtp2:e4facts:big_in_casc_mean}~ms inside the cascade against
\N{chk:mtp2:e4facts:big_alone_esc_mean}~ms for the same items in the qwen-alone run an hour earlier;
Section~\ref{sec:ablation} tests whether prefix caching explains it. The largest error, +42\% on E4
entities, came from the session's very first run: Nox's standalone E4-entity pass averaged
\N{chk:mtp2:e4ent:small_mean}~ms with a 95th percentile of \N{chk:mtp2:e4ent:small_p95}~ms, against
\N{chk:mtp2:e4ent:small_in_casc_mean}~ms inside the later cascade run. A replay is only as good as
its service-time sample.

\subsection{Throughput: a two-station model is close, and the naive one is not}
Table~\ref{tab:tput} compares the measured cascade throughput at $c=8$ with the three predictions of
Section~\ref{sec:replay}. Ignoring contention overpredicts throughput 3--5$\times$. Both the bottleneck
bound and the load-dependent two-station model are within 3\% when one station dominates (E4 entities,
where no item escalates, and E2, where nearly all do), and overpredict by 14--25\% when the two are
comparably loaded (E3, E4 facts). That is where such models are known to be loosest. The decision-model
runtime is nearly serial (Nox alone gains only 8--43\% throughput from 8 clients, against about
2$\times$ for qwen), so at 33--40\% escalation Nox and qwen saturate at similar rates. Clients
then wait alternately on each stage, which a model built from per-station averages misses. The
bottleneck bound still identifies the right design lever: on these tasks a faster or replicated
first stage, not a faster qwen, is what would raise cascade throughput further.

\begin{table}[t]
\centering
\small
\begin{tabular}{lrrrrrrrr}
\toprule
 & & \multicolumn{4}{c}{Cascade throughput, $c=8$ (items/s)} & \multicolumn{3}{c}{Prediction error} \\
\cmidrule(lr){3-6}\cmidrule(lr){7-9}
Task & Esc. & no contention & bound & model & measured & no cont. & bound & model \\
\midrule
E3 triage & 33\% & 30.4 & 12.31 & 11.86 & 9.84 & +210\% & +25\% & +21\% \\
E4 fact typing & 40\% & 27.3 & 10.22 & 10.02 & 8.79 & +210\% & +16\% & +14\% \\
E4 entity typing & 0\% & 40.5 & 7.92 & 8.10 & 7.97 & +408\% & -1\% & +2\% \\
E2 spelling & 98\% & 6.0 & 1.77 & 1.75 & 1.72 & +247\% & +3\% & +1\% \\
\bottomrule
\end{tabular}

\caption{Cascade throughput at 8 concurrent clients: three analytical predictions against the
measurement. ``No contention'' is $c/\mathbb{E}[\hat T]$ with single-stream service times; ``bound'' is
$\min(X_S^{\max}, X_B^{\max}/p)$ from the standalone 8-client runs; ``model'' is the load-dependent
two-station closed network of Section~\ref{sec:replay}. Errors are relative to the measurement.}
\label{tab:tput}
\end{table}

\subsection{Serving configuration and reproducibility}
\label{sec:ablation}
The deployed qwen uses two features that plausibly break assumptions A1 and A2: MTP speculative
decoding, which verifies several drafted tokens in one forward pass and so changes the numerical
path of a greedy decode; and automatic prefix caching (APC), which on this hybrid
attention/linear-attention model restores recurrent state from cached blocks instead of recomputing
it. We restarted qwen twice, with MTP off and then with MTP and APC both off, and repeated the qwen
arm and the live cascade at both concurrency levels. The decision model and the gates were unchanged.
Each configuration also got a second single-stream qwen pass over E3, identical to the first,
except that with APC on every prompt is now a cache hit. Table~\ref{tab:ablation} summarises.

\begin{table}[t]
\centering
\small
\setlength{\tabcolsep}{4pt}
\resizebox{\textwidth}{!}{\begin{tabular}{llrrrrrrrrr}
\toprule
 & & \multicolumn{4}{c}{qwen alone} & \multicolumn{2}{c}{$\Delta$ vs qwen (pt)} & \multicolumn{2}{c}{Speed-up $c=1$} & Gain \\
\cmidrule(lr){3-6}\cmidrule(lr){7-8}\cmidrule(lr){9-10}
qwen config & Task & ms & repro & $c1{=}c8$ & $c1{=}c1$ & pred. & meas. & pred. & meas. & $c=8$ \\
\midrule
\multirow{3}{*}{\shortstack[l]{MTP-2,\\prefix cache on (deployed)}} & E3 & 528 & 99.0 & 99.1 & 100.0 & +0.0 & +0.1 & 2.01 & 2.19 & 2.41 \\
 & E4 facts & 502 & 96.9 & 97.0 & -- & -0.1 & +0.0 & 1.71 & 1.80 & 2.16 \\
 & E2 & 1077 & 93.2 & 93.4 & -- & -1.2 & -1.2 & 0.81 & 0.88 & 0.99 \\
\midrule
\multirow{3}{*}{\shortstack[l]{MTP off,\\prefix cache on}} & E3 & 605 & 99.0 & 99.5 & 99.8 & +0.0 & +0.0 & 2.09 & 1.58 & 2.71 \\
 & E4 facts & 570 & 96.6 & 97.5 & -- & -0.1 & +0.0 & 1.78 & 1.60 & 2.48 \\
 & E2 & 1736 & 94.4 & 94.4 & -- & -1.2 & -2.0 & 0.87 & 0.93 & 1.01 \\
\midrule
\multirow{3}{*}{\shortstack[l]{MTP off,\\prefix cache off}} & E3 & 611 & 98.8 & 99.0 & 100.0 & +0.0 & +0.1 & 2.09 & 2.26 & 2.50 \\
 & E4 facts & 574 & 97.6 & 97.0 & -- & -0.1 & -0.1 & 1.78 & 1.86 & 2.06 \\
 & E2 & 1392 & 93.4 & 94.6 & -- & -1.2 & -1.2 & 0.85 & 0.92 & 1.02 \\
\bottomrule
\end{tabular}
}
\caption{The live check under three qwen serving configurations. \emph{qwen alone}: mean
single-stream latency (ms); share of items on which it reproduces its frozen answer (``repro'');
agreement of its own single-stream and 8-client answers in the same session ($c1{=}c8$), and of two
identical single-stream passes ($c1{=}c1$, E3 only). $\Delta$: cascade minus qwen, predicted and
measured. Speed-up: qwen-alone over cascade mean latency at $c=1$, predicted by the replay (with
that configuration's service times) and measured. Gain: cascade over qwen-alone throughput at $c=8$.}
\label{tab:ablation}
\end{table}

\paragraph{Answers: deterministic per configuration and load, not across them.} In the deployed
configuration two identical single-stream passes agreed on all 864 E3 items
(\N{chk:mtp2:e3:big1_rep_agree}\%), so qwen is repeatable when nothing changes. What changes its
answers is the serving \emph{context}. Single-stream and 8-client answers in the same session differ
on 0.5--3\% of E3 and E4-facts items and 5--7\% of E2 items in every configuration, consistent with
batch composition changing the reduction order of the kernels. The configuration itself also matters: single-stream answers with MTP on and off
agree on \N{xcfg:e3:mtp2:mtp0}\% of E3 and \N{xcfg:e4facts:mtp2:mtp0}\% of E4 facts, and on only
\N{xcfg:e2:mtp2:mtp0}\% of E2. Turning MTP off did not make qwen reproduce its frozen answers any
better (E4 facts \N{chk:mtp0:e4facts:big1_agree}\% vs \N{chk:mtp2:e4facts:big1_agree}\%), nor did
turning APC off as well (\N{chk:mtp0apc0:e4facts:big1_agree}\%). The frozen answers were produced
under one particular configuration and load (MTP-2, APC, 8-way), and no single-stream configuration
recovers them. MTP is therefore not the cause of the drift; it is one of several things the answers
depend on. In every configuration the replay's predicted cascade-minus-qwen difference stayed within
0.1 points of the measurement on E3 and E4 facts.

\paragraph{Latency: prefix-cache hits are a hidden variable.} With MTP off and APC on, the second,
fully cached pass over E3 was \emph{slower} than the first: \N{chk:mtp0:e3:big1_rep_mean}~ms
against \N{chk:mtp0:e3:big_mean}~ms (+\N{chk:mtp0:e3:rep_slowdown}\%), and it changed 2 of 864
answers. In the deployed MTP-2 configuration the cached pass was instead slightly faster
(\N{chk:mtp2:e3:big1_rep_mean} vs \N{chk:mtp2:e3:big_mean}~ms) with identical answers. The serving image carries a two-line fix for a prefix-cache
bug on this model, in which a cache hit restored an all-zero recurrent state because the state slot was
computed from the smallest KV-group block size (8 tokens at MTP-2) instead of the state block size.
That fix was validated as bit-identical between cold and cached outputs \emph{at MTP-2}. Our MTP-off
result, where cache hits changed answers and ran 38\% slower, is consistent with the fix not covering
the block sizes that arise with MTP off. We did not investigate further. For a replay it means a
serving flag nobody thinks of as affecting answers can change both answers and service times. It hit the cascade directly: every escalated prompt had just been
sent by the qwen-alone run, so inside the MTP-off cascade qwen took
\N{chk:mtp0:e3:big_in_casc_mean}~ms per escalated E3 item against
\N{chk:mtp0:e3:big_alone_esc_mean}~ms for the same items alone. The replay, built from cold-cache
service times, underpredicted cascade latency by a quarter (\N{chk:mtp0:e3:comp_mean} against
\N{chk:mtp0:e3:meas_mean}~ms) and overpredicted its speed-up (\N{chk:mtp0:e3:speedup_pred}$\times$
vs \N{chk:mtp0:e3:speedup_meas}$\times$). With APC off the effect disappears and the replay is again
within 5--8\% (\N{chk:mtp0apc0:e3:comp_mean} vs \N{chk:mtp0apc0:e3:meas_mean}~ms on E3).

The small residual in the other direction, qwen answering escalated items 3--7\% faster inside the
cascade than in a back-to-back qwen-alone run, persists with APC off
(\N{chk:mtp0apc0:e3:big_in_casc_mean} vs \N{chk:mtp0apc0:e3:big_alone_esc_mean}~ms), so it is not
caching. It coincides with the cascade's spacing of qwen requests (each is preceded by a Nox call
and often by several non-escalated items). We did not isolate it further.

\paragraph{MTP is worth keeping for this workload.} Even though every answer here is a short JSON
object, MTP-2 cut qwen's single-stream latency from \N{chk:mtp0apc0:e3:big_mean} to
\N{chk:mtp2:e3:big_mean}~ms on E3 and from \N{chk:mtp0apc0:e2:big_mean} to
\N{chk:mtp2:e2:big_mean}~ms on E2, whose answers are longer. The deployed configuration is the
fastest and the best-behaved of the three for the cascade.

\section{Discussion}

\paragraph{What the replay is good for.} It answered the design question correctly and cheaply. Of
twenty (first stage, task) pairs it singled out one pair to deploy on two tasks, one task needing no
cascade and one task to leave alone, and every one of those calls survived the live test. It also
explained \emph{why} the other first stages fail (too weak, so everything escalates; or too slow, so
nothing is saved), which a single live run of one configuration would not have shown. Its
accuracy-difference predictions were within 0.1 points. Exploring the space live would have meant
running each candidate decision model on the GB10, re-serving it for each experiment, and measuring
five cascades on four tasks under two load levels. In the replay that is a loop over saved records.

\paragraph{What it silently assumes.} Three assumptions did the work, and each was violated
somewhere:
\begin{enumerate}
  \item \emph{The gate's inputs are reproducible.} They were, exactly, because a decision model's head
        is a deterministic function of the input. This is the property that makes the cascade's
        routing predictable, and it is a reason to put a decision model, rather than a sampled
        generative judge, at the gate.
  \item \emph{The expensive model's answers are reproducible.} They were not, at the 1--7\% level, even
        at temperature~0. The replay absorbed this because both arms of the comparison inherit the same
        model, but a replay that reports absolute accuracy should carry the run-to-run spread of the
        generative model as an error bar. The ablation shows this drift is not one feature's fault: no
single-stream configuration we tried, with or without speculative decoding and prefix caching,
reproduced the frozen answers better than another. The answers are a function of model, serving
configuration \emph{and} load.
  \item \emph{Service times are measured in the regime being predicted.} Mixing an 8-way concurrent
        sample with single-stream ones produced a 3--4$\times$ error that looked plausible: it was
        the latency we had first quoted when proposing the cascade. Throughput under load needs a
        queueing model, and the simplest one that respects per-station saturation was within 25\%,
        while ``no contention'' was off by 3--5$\times$.
\end{enumerate}

\paragraph{Practical recipe.} Freeze the evaluation inputs; run every candidate once on all of them,
recording answer, confidence and request time, with each model's service time measured in the load
regime it will serve in; replay the cascade over the records with held-out thresholds targeting the
objective you actually have (here parity, with the expensive error class constrained separately);
then validate the chosen point live, with the gates frozen beforehand, and compare item by item.
The item-level comparison is what separates ``the prediction was right'' from ``the errors
happened to cancel''.

\paragraph{Local hardware.} Two observations are specific to running both stages on DGX Sparks.
The decision-model runtime is nearly serial on the GB10, so under load the 4B first stage, not the
generative model, sets the cascade's throughput ceiling; replicating it is cheap in memory
($\sim$19~GB) and is the next lever. The cascade also turns two boxes into a pipeline: every item the
decision model settles on spark2 in about 90~ms is a half-second generative request that spark1
never queues.

\section{Limitations}
\begin{itemize}
  \item \textbf{One live pair.} Only Nox-4B $\rightarrow$ qwen was run live; the other four first
        stages are replay-only. Their accuracy predictions rest on the same determinism that held for
        Nox (all are deterministic-head decision models from the companion study), but their latency
        predictions use the companion study's service times and were not re-measured.
  \item \textbf{Evaluated on the dataset the gates were fitted on.} Thresholds were fitted 2-fold by
        session, so no item's label set its own threshold, but the live run used the same frozen items
        as the replay. That is deliberate (it isolates prediction error from sampling error), but it
        is not a test on new data. The two folds' thresholds differ (0.70 vs 0.72 on E4 facts, 0.83
        vs 0.88 on E3), and a deployment would need its own calibration set, as System One Auto
        notes~\cite{sysoneauto}.
  \item \textbf{One hardware and software profile.} Two GB10 systems on one LAN, one vLLM build, one
        decision-model runtime. The serial behaviour of that runtime under concurrency shapes the
        throughput results and may not hold for other runtimes.
  \item \textbf{Small label sets for some tasks.} E4 entity typing has 92 items; its CIs span
        $\pm$5 points.
  \item \textbf{The E3 metric is ours.} E3 is scored here as binary accuracy with a missed-name
        constraint, not with the companion study's forwarded-share-at-$k$-misses metric.
  \item \textbf{Non-reasoning qwen only.} The thinking variant, which is the only candidate that
        could raise the accuracy ceiling, was not tested as a second stage.
\end{itemize}

\section{Conclusion}
Replaying frozen model outputs is a cheap and, within its assumptions, accurate way to design a
cascade. Here it chose a two-box decision-model/generative-model cascade that kept the generative
model's accuracy at roughly twice its speed. It also predicted that accuracy difference within a tenth
of a point and routed every item exactly as the live system did. The assumptions are what a
practitioner has to check: a deterministic gate (decision models provide one), a second stage whose
answer drift is measured rather than assumed away, and service times sampled in the load regime and
cache state being predicted. Get the last one wrong and a replay that looks authoritative is off by a
factor of three. Measure them properly and a few seconds of arithmetic stands in for an afternoon of
GPU time.

\section*{Data and code availability}
The replay model (\texttt{clef/eval/cascade.py}), the live router (\texttt{cascade\_live.py}), the
script that generates every number, table and figure in this paper from the result files
(\texttt{cascade\_paper.py}) and the frozen dataset's manifest are in the public
repository~\cite{dgxfun}. The campaign data itself is not released because it quotes play sessions
verbatim and names real people.

\section*{Use of AI assistance}
The experiments, analysis and a first draft of this paper were carried out by an AI coding agent
(Claude, Anthropic) working under the author's direction, with the author reviewing results and
making the decisions recorded in the paper. The author takes full responsibility for the content.

\bibliographystyle{plainnat}
\begin{thebibliography}{10}

\bibitem[Roussos(2026a)]{companion}
K.~Roussos.
\newblock Decision models at the game table: where non-generative classifiers help and fail in LLM
  tooling for tabletop role-playing campaigns.
\newblock Manuscript, October 2026. In~\cite{dgxfun}, \texttt{paper/main.tex}.

\bibitem[TypeSafe(2026)]{typesafe}
TypeSafe.
\newblock Jev and the System One API.
\newblock Documentation, \url{https://docs.typesafe.ai}, accessed October 2026.

\bibitem[vLLM Semantic Router Team(2026)]{sysoneauto}
vLLM Semantic Router Team.
\newblock System One Auto.
\newblock Blog post, \url{https://vllm-sr.ai/blog/system-one-auto/}, accessed October 2026.

\bibitem[Chen et~al.(2023)]{frugalgpt}
L.~Chen, M.~Zaharia, and J.~Zou.
\newblock FrugalGPT: How to use large language models while reducing cost and improving performance.
\newblock arXiv:2305.05176, 2023.

\bibitem[Ong et~al.(2024)]{routellm}
I.~Ong, A.~Almahairi, V.~Wu, W.-L. Chiang, T.~Wu, J.~E. Gonzalez, M.~W. Kadous, and I.~Stoica.
\newblock RouteLLM: Learning to route LLMs with preference data.
\newblock arXiv:2406.18665, 2024.

\bibitem[Little(1961)]{little1961}
J.~D.~C. Little.
\newblock A proof for the queuing formula: $L = \lambda W$.
\newblock \emph{Operations Research}, 9(3):383--387, 1961.

\bibitem[Reiser and Lavenberg(1980)]{reiser1980}
M.~Reiser and S.~S. Lavenberg.
\newblock Mean-value analysis of closed multichain queuing networks.
\newblock \emph{Journal of the ACM}, 27(2):313--322, 1980.

\bibitem[Roussos(2026b)]{dgxfun}
K.~Roussos.
\newblock dgx-fun: experiments with local LLM serving on NVIDIA DGX Spark.
\newblock \url{https://github.com/kostadis/dgx-fun}, 2026.

\end{thebibliography}

\end{document}
